\documentclass[11pt,a4paper,reqno]{amsart}

\usepackage[a4paper,
            textwidth=15.5cm,
            textheight=23cm,
            hcentering,
            vcentering,
            headheight=14pt,
            headsep=0.6cm,
            footskip=1.1cm]{geometry}

\usepackage[T1]{fontenc}
\usepackage[utf8]{inputenc}
\usepackage{amsmath,amsthm,amscd}
\usepackage{mathtools}
\usepackage[scaled=0.97]{newpxtext}   % Palatino text
\usepackage[scaled=0.97]{newpxmath}   % matching Palatino math (includes AMS symbols)
\usepackage{microtype}
\usepackage{stmaryrd}                 % \llbracket, \rrbracket

%% --- Utilities ---------------------------------------------------
\usepackage{graphicx}
\usepackage{caption}
\usepackage{verbatim}
\usepackage{enumitem}
\usepackage{xcolor}
\definecolor{linkblue}{RGB}{0,51,128}
\usepackage[colorlinks=true,
            linkcolor=linkblue,
            citecolor=linkblue,
            urlcolor=linkblue,
            pdfusetitle,
            bookmarksnumbered]{hyperref}
\usepackage[nameinlink,capitalize,noabbrev]{cleveref}

%% --- Running heads and spacing ----------------------------------
\pagestyle{headings}
\setlength{\parskip}{2pt plus 1pt}
\allowdisplaybreaks[1]
\hyphenation{semi-stable Hil-bert Fro-be-nius Castel-nuovo Mum-ford}

%% --- Theorem environments ----------------------------------------
\numberwithin{equation}{section}

\theoremstyle{plain}
\newtheorem{theorem}{Theorem}[section]
\newtheorem{lemma}[theorem]{Lemma}
\newtheorem{proposition}[theorem]{Proposition}
\newtheorem{corollary}[theorem]{Corollary}
\newtheorem{conjecture}[theorem]{Conjecture}
\newtheorem*{theoremA}{Theorem A}
\newtheorem*{theoremB}{Theorem B}
\newtheorem*{theoremC}{Theorem C}
\newtheorem*{theoremD}{Theorem D}

\theoremstyle{definition}
\newtheorem{definition}[theorem]{Definition}
\newtheorem{example}[theorem]{Example}

\theoremstyle{remark}
\newtheorem{remark}[theorem]{Remark}
\newtheorem*{acknowledgement}{Acknowledgement}

\setlist[enumerate]{label=(\arabic*),leftmargin=2em,itemsep=2pt,topsep=5pt}

%% --- Macros ------------------------------------------------------
%% Standard notation used throughout:
%%   \reg_R M          Castelnuovo--Mumford regularity
%%   I^{[q]}           Frobenius power, q=p^e
%%   \mm               the homogeneous maximal ideal (x,y,z,u,w)R
%%   A_q               k[x,y,z]/(x^q,y^q,z^q);  (A_q)_{d,i} degree d, weight i mod 3
%%   b_q, c_q          lower bound and coefficient degree of the primitive kernel vector
%%   \length           length of a module (typeset as \ell)
%%   \llbracket x\rrbracket   power series
\newcommand{\mm}{\mathfrak m}
\newcommand{\kk}{k}
\newcommand{\Id}{\mathrm{Id}}
\newcommand{\length}{\ell}

\DeclareMathOperator{\Ann}{Ann}
\DeclareMathOperator{\Ass}{Ass}
\DeclareMathOperator{\charac}{char}
\DeclareMathOperator{\coker}{coker}
\DeclareMathOperator{\depth}{depth}
\DeclareMathOperator{\diag}{diag}
\DeclareMathOperator{\End}{end}
\DeclareMathOperator{\Ext}{Ext}
\DeclareMathOperator{\gr}{gr}
\DeclareMathOperator{\Hom}{Hom}
\DeclareMathOperator{\ini}{in}
\DeclareMathOperator{\NF}{NF}
\DeclareMathOperator{\ord}{ord}
\DeclareMathOperator{\Proj}{Proj}
\DeclareMathOperator{\reg}{reg}
\DeclareMathOperator{\Spec}{Spec}
\DeclareMathOperator{\supp}{supp}
\DeclareMathOperator{\Syz}{Syz}
\DeclareMathOperator{\wt}{wt}

%%%%%%%%%%%%%%%%%%%%%%%%%%%%%%%%%%%%%%%%%%%%%%%%%%%%%%%%%%%%%%%%%%%%%%%%%%%%%%%%%%%

\begin{document}

%%%%%%%%%%%%%%%%%%%%%%%%%%%%%%%%%%%%%%%%%%%%%%%%%%%%%%%%%%%%%%%%%%%%%%%%%%%%%%%%%%%

\title{Superlinear regularity of Frobenius powers on a normal hypersurface}

\author[T.N.Trung]{Tran Nam Trung}
\address{Institute for Artificial Intelligence, VNU University of Engineering and Technology, Hanoi, Vietnam}
\email{tntrung@vnu.edu.vn}

\author[H.L.Truong]{Hoang Le Truong}
\address{Institute for Artificial Intelligence, VNU University of Engineering and Technology, Hanoi, Vietnam}
\email{hltruong@vnu.edu.vn, truonghoangle@gmail.com}


\thanks{ChatGPT (OpenAI) was used in an iterative drafting process based on the authors' research notes and detailed instructions, and to assist with subsequent revisions of the manuscript. The authors reviewed and substantially revised the resulting text, verified the mathematical arguments and references, and take responsibility for the final manuscript.}


\keywords{Frobenius powers, Castelnuovo--Mumford regularity, Local cohomology, Primary decomposition, Associated primes, Gr\"obner bases}

\subjclass[2020]{13A35, 13D45, 13C13, 13P10}

\begin{abstract}
We give answers to four questions on linear bounds for Frobenius powers.
Let $\mathfrak m$ denote the homogeneous maximal ideal.
We construct a height-two prime $I$ in a geometrically normal four-dimensional
graded hypersurface $R$ of characteristic two such that, for every power
$q\geq16$ of four, the Loewy length of $H^0_{\mathfrak m}(R/I^{[q]})$
is at least $q^{3/2}/8-1$. Thus no fixed integer $C$ gives annihilation
by $\mathfrak m^{Cq}$ for all $q$, as required by the Hochster--Huneke
condition $(\mathrm{LC})$. The same construction gives superlinear lower
bounds for the Castelnuovo--Mumford regularity in Katzman's question,
for the exponents of the embedded primary components in the question of
Smith and Swanson, and for the degree in the least variable in Katzman's
Gr\"obner basis conjecture. The bounds follow from the degree of the kernel
line bundle of a multiplication pencil associated to a family of plane quartics.
\end{abstract}

\maketitle

%%%%%%%%%%%%%%%%%%%%%%%%%%%%%%%%%%%%%%%%%%%%%%%%%%%%%%%%%%%%%%%%%%%%%%%%%%%%%%

\section{Introduction}

Let $R$ be a standard graded algebra over a field of characteristic
$p>0$, let $\mm=R_+$, and let $I=(f_1,\ldots,f_s)$ be homogeneous.
For $q=p^e$, the ideal $I^{[q]}=(f_1^q,\ldots,f_s^q)$ defines the
scheme-theoretic inverse image of $V(I)$ under the $e$th iterate of
the absolute Frobenius morphism of $\Spec R$. Its support is independent
of $e$, and its displayed generators have degrees $q\deg f_1,\ldots,q\deg f_s$.
The embedded structure at
the vertex can nevertheless grow faster. We study four linear bounds
that this phenomenon obstructs.

Throughout the paper, let $k$ be a field of characteristic two and let
$P=k[x,y,z,u,w]$ with its standard grading. Set
\begin{equation}\label{eq:fixed-pair}
\begin{gathered}
 h=(x^2+yz)(y^2+xz),\quad J=(uz^4+wh),\\
 K=(x,y,z)P,\qquad R=P/J,\qquad I=KR .
\end{gathered}
\end{equation}
For $q=2^e\geq4$ put $b_q=2^{e+\lfloor e/2\rfloor-3}-1$.
For $q=4^r$ this gives $b_q=q^{3/2}/8-1$.

The pencil of plane quartics defined by $uz^4+wh$ relates the geometry
over $[u:w]$ to the embedded structure at the vertex. The orders of its
maximal minors determine the degree of a kernel line bundle on
$\mathbb P^1$. Graded local duality recovers from this degree a summand
of the local cohomology supported at the vertex.

\subsection*{Local-cohomology annihilation}
The quotient
\[
 L_q=H^0_{\mm}(R/I^{[q]})
     =(I^{[q]}:\mm^\infty)/I^{[q]}
\]
measures the defect of saturation at the vertex. The condition
$(\mathrm{LC})$ asks for a constant $B$ such that
$\mm^{Bq}L_q=0$ for every $q$. It arose in the study of tight closure
by Hochster and Huneke~\cite[discussion after Proposition 4.16]{HH}.
Vraciu~\cite{Vraciu} studied this annihilation exponent for graded
affine algebras. Huneke~\cite[Corollary 3.2]{HunekeLC} proved
$(\mathrm{LC})$ for homogeneous ideals with one-dimensional quotient
in equidimensional graded rings. Asgharzadeh~\cite[Theorem 1.2(i)]{Asgharzadeh}
proved that every homogeneous ideal in a normal standard graded
Cohen--Macaulay algebra of dimension at most three satisfies
$(\mathrm{LC})$.
Part (ii) of the same theorem uses regularity of the ambient ring at the
relevant primes on the support. Here the generic point $I$ of the plane
has $\dim R/I=2$, and $R_I$ has dimension two and embedding dimension
three since $uz^4+wh\in(x,y,z)^4P$. See also \cite[Conjecture 3.1]{PQ} for the local formulation.

\subsection*{Katzman's regularity question}
For fixed homogeneous ideals $J,K$ in a polynomial ring $P$,
Katzman~\cite[Question 2]{Katzman} and Katzman--Zhang
\cite[Question 1.4]{KZ} ask whether $\reg_P(J+K^{[q]})=O(q)$.
When $J=0$, Frobenius is flat and multiplies the shifts in a graded
free resolution by $q$. Brenner~\cite{Brenner} proved linear bounds
under geometric hypotheses. Chardin~\cite{Chardin} studied regularity
under Frobenius using local cohomology of graded complexes and discussed
Katzman's question in \cite{ChardinSurvey}. Linear bounds are known for
graded rings of finite Frobenius representation type
\cite[\S7]{KSSZ}. Bounds in terms of $F$-thresholds and $F$-pure thresholds are given in \cite{DSNP}.
Katzman and Zhang related bounds for a principal fixed ideal $J$ to
discreteness of $F$-jumping coefficients~\cite[Theorem 3.3]{KZ}.
Our example has $J$ principal. Throughout, regularity means
$\reg_R M=\max_i\{\End H^i_{\mm}(M)+i\}$.

\subsection*{Primary decomposition}
Ordinary powers of an ideal in a Noetherian ring admit primary
decompositions with linear bounds on the powers of the radicals
\cite{Swanson}. Yao~\cite{Yao} proved the corresponding bound for
products of ordinary powers. For Frobenius powers, Smith and Swanson
\cite[p.~3071]{SS} ask whether one can choose a fixed integer $\ell$
and, for each $q$, a primary decomposition satisfying
\begin{equation}\label{eq:SS-question}
 I^{[q]}=\bigcap_i Q_{i,q},\qquad
 (P_{i,q}^{[q]})^\ell\subseteq Q_{i,q},
 \quad P_{i,q}=\sqrt{Q_{i,q}}.
\end{equation}
Sharp and Nossem~\cite{SN} relate this condition to localization
of tight closure. Dinh~\cite[Theorem 2.2]{Dinh} proved it for
homogeneous ideals with one-dimensional quotient in positively
graded algebras over a field. Dinh also exhibited linear growth for
Frobenius powers whose associated primes have infinite union~\cite{Dinh}.
Here $R/I$ has dimension two. In every reduced decomposition the exponent
of the component at the vertex grows superlinearly.

\subsection*{The least variable in a Gr\"obner basis}
Fix graded reverse lexicographic order on $P$ and let $w$ be its least
variable. Katzman~\cite[Conjecture 4]{Katzman} asks whether the
largest $w$-degree in the reduced Gr\"obner basis of $J+K^{[q]}$
is $O(q)$. This question arose from a criterion for localization of
tight closure at one element. Hermiller and Swanson~\cite[Theorem 2.1]{HS}
proved the bound for arbitrary ideals in one or two variables. A bound on the $w$-degree gives a bound on the exponent needed to saturate with respect to $w$. We exhibit a class on which a superlinear power of $w$
still acts nontrivially.

\begin{theoremA}[Theorems~\ref{thm:regularity}, \ref{thm:LC}, \ref{thm:primary} and~\ref{thm:groebner}]
The ring $R$ is a geometrically normal four-dimensional hypersurface
and $I$ is a height-two prime. For each $q=2^e\geq8$ the following hold.
\begin{enumerate}
\item $w^{b_q-1}H^0_{\mm}(R/I^{[q]})\ne0$.
\item Both $\reg_R I^{[q]}$ and $\reg_P(J+K^{[q]})$ are at least
$3q/2+b_q+1$.
\item Every reduced primary decomposition of $I^{[q]}$ has an
$\mm$-primary component $Q$ with $w^{b_q-1}\notin Q$.
\item For graded reverse lexicographic order with $w$ least, the
reduced Gr\"obner basis of $J+K^{[q]}$ contains an element of
$w$-degree at least $b_q$.
\end{enumerate}
\end{theoremA}

The first, second and fourth bounds are superlinear because
$b_q/q\to\infty$. For the third, a containment
$(\mm^{[q]})^\ell\subseteq Q$ would give $w^{q\ell}\in Q$, hence
$q\ell\geq b_q$. Thus no fixed $\ell$ satisfies~\eqref{eq:SS-question}.
Corollary~\ref{cor:dimension} gives the same superlinear lower bound on Loewy length in every dimension at least four in characteristic two.

\subsection*{Geometry of the proof}
The plane $Z=V(I)$ lies in the singular locus of $\Spec R$. Its $e$th Frobenius inverse image is $Z_e=\Spec(R/I^{[q]})$, with $q=2^e$.
The projectivized normal cone along $Z$ has a plane quartic as its
fibre over a general point $(u,w)$. Its equation depends only on
$[u:w]$. On $w\ne0$ it is $h+\alpha z^4=0$, with $\alpha=u/w$.
These are Monsky's line-$S_4$ quartics~\cite{Monsky}, also used by
Borevitz et al.~\cite{Borevitz} to study localization of tight closure
through the escape dynamics of the parameter. 

In a fixed degree and character, multiplication by $uz^4+wh$ gives
a pencil over $k[u,w]$ whose kernel on the parameter line is
$\mathcal O_{\mathbb P^1}(-c_q-1)$. This defines the kernel index $c_q$. Monsky's fibre lengths locate the finite nonzero parameters
at which the rank drops. To compute the kernel degree, we also need
the orders of the maximal minors at those parameters. We obtain them
by lifting a syzygy of the smooth quartic along $h+(\alpha+t)z^4$.
At infinity the fibre is the quadruple line $z^4=0$, where division
by a monic quartic gives the remaining order. An upper bound for this
order yields $c_q\geq b_q$, which suffices for Theorem A. Its exact
value gives $c_q=b_q$.

Graded local duality then identifies the dual of a summand of $L_q$,
up to shift, with $k[u,w]/(u,w)^{c_q}$. The highest nonzero degree of the summand gives the regularity bound.
The action of $w$ on it gives the other three bounds.
This summand is supported at the affine vertex. Its sheafification
on $\Proj R$ is zero. The same family also determines all the associated
primes of $R/I^{[q]}$ (Theorem~\ref{thm:ass-exact}): the height-one
primes over $k[u,w]$ come from jumps in fibre length, and $L_q$
detects the vertex.

\subsection*{Conventions}
Throughout, $k$ is a field of characteristic two, $q=2^e$ runs
through the powers of two, and a dyadic integer means a power of
two. The ring $P=k[x,y,z,u,w]$, the ideals $J,K$, the hypersurface
$R=P/J$ and the prime $I=KR$ are those of~\eqref{eq:fixed-pair};
$\mm=(x,y,z,u,w)R$ is the homogeneous maximal ideal and
$L_q=H^0_{\mm}(R/I^{[q]})$. For a graded module $M$ we write
$\End M$ for the largest degree in which $M$ is nonzero and
$\reg_R M=\max_i\{\End H^i_{\mm}(M)+i\}$ for its Castelnuovo--Mumford
regularity. The length of a module $M$ is $\length(M)$. We write
$A_q=k[x,y,z]/(x^q,y^q,z^q)$, and $(A_q)_{d,i}$ for its component of
degree $d$ and weight $i$ modulo three for the auxiliary grading
introduced in Section~\ref{sec:pencil}; the words ``weight'' and
``character'' are used interchangeably for this grading. The integer
$b_q=2^{e+\lfloor e/2\rfloor-3}-1$ is the lower bound of Theorem A,
and $c_q$ is the coefficient degree of the primitive kernel vector
defined in Section~\ref{sec:pencil}. For a ring $A$ we write
$A\llbracket t\rrbracket$ for the power series ring, and $\ini(g)$
for the initial term of a polynomial $g$ with respect to a fixed
monomial order.

%%%%%%%%%%%%%%%%%%%%%%%%%%%%%%%%%%%%%%%%%%%%%%%%%%%%%%%%%%%%%%%%%%%%%%%%%%%%%%

\section{The central pencil and local cohomology}\label{sec:pencil}
Throughout, $k$ has characteristic two and $q$ is a power of two.
Field extension commutes with the matrix calculations and local
cohomology, so we may extend $k$ when computing ranks and local orders.

\subsection{The hypersurface and its quartic family}\label{subsec:transverse}

\begin{proposition}\label{prop:normal}
The ring $R$ in~\eqref{eq:fixed-pair} is a four-dimensional geometrically
normal hypersurface domain, and $I$ is a height-two prime.
\end{proposition}
\begin{proof}
Since $\gcd(z^4,h)=1$, Gauss's lemma shows that $uz^4+wh$ is irreducible.
Thus $R$ is a four-dimensional Cohen--Macaulay domain. Over an algebraic
closure, the radical of its Jacobian ideal is
\[
(z,xy,wx,wy)=(z,x,y)\cap(z,y,w)\cap(z,x,w).
\]
Indeed the partial derivatives are $wz(x^2+yz)$, $wz(y^2+xz)$,
$w(x^3+y^3)$, $z^4$, and $h$. The singular locus has dimension two.
Serre's criterion gives normality, and the same calculation after field
extension proves geometric normality. Finally $R/I=k[u,w]$.
\end{proof}

The Jacobian calculation shows that the plane $Z=V(I)$ lies in the
singular locus of $X=\Spec R$. Its normal cone explains the quartics
used below. Give $x,y,z$ degree one and $u,w$ degree zero. The equation
of $X$ is homogeneous of degree four, so
\begin{equation}\label{eq:normal-cone}
\gr_I R\simeq
\frac{k[u,w,\xi,\eta,\zeta]}
 {u\zeta^4+w(\xi^2+\eta\zeta)(\eta^2+\xi\zeta)}.
\end{equation}
Here $\xi,\eta,\zeta$ are the initial forms of $x,y,z$.
Thus the fibre of $\Proj_Z(\gr_I R)$ over a point
$(u,w)=(a,b)$ with $ab\ne0$ is the smooth plane quartic
$C_{a/b}=V(h+(a/b)z^4)$. Smoothness is checked in
Section~\ref{subsec:quartic-syzygies}.
Equivalently, at the generic point of $Z$, with $K=k(u,w)$,
\[
\widehat{R_I}\simeq K\llbracket x,y,z\rrbracket/(h+(u/w)z^4).
\]
This identifies the transverse surface singularity with the completed
cone over the generic quartic. 
\subsection{The central multiplication pencil}
Put $S=k[x,y,z]$, $S'=k[u,w]$, and $T_q=R/I^{[q]}$.
For $q\geq4$, set
\[
A_q=S/(x^q,y^q,z^q),\qquad d_q=3q/2-3,\qquad n_q=q^2/4-2.
\]
In addition to ordinary degree, give $x,y,z$ weights $1,2,0$ in
$\mathbb Z/3\mathbb Z$. Write $(A_q)_{d,i}$ for the part of degree $d$
and weight $i$. We use ``weight'' and ``character'' for this auxiliary
grading. Both $z^4$ and $h$ have weight zero.
The central pencil is
\begin{equation}\label{eq:central-pencil}
M_q(u,w)=u\mu_{z^4}+w\mu_h:
S'(-1)\otimes_k(A_q)_{d_q,1}
\longrightarrow S'\otimes_k(A_q)_{d_q+4,1}.
\end{equation}
We write $\mathcal N_q=\coker M_q$.
At a finite parameter $\alpha=u/w$, the pencil is multiplication by
$h+\alpha z^4$. The following proposition gives its generic rank and
also excludes $\alpha=0$ from its rank-drop locus.

\begin{proposition}\label{prop:rank}
For every dyadic $q\geq4$, multiplication by $h$ maps
$(A_q)_{3q/2-3}$ onto $(A_q)_{3q/2+1}$.
On character one, the source and target have dimensions
$q^2/4-1$ and $q^2/4-2$, respectively.
\end{proposition}
\begin{proof}
We may assume that $k$ is algebraically closed. Write $q=2Q$ and
$j=3Q+1$. Put $a=x^2+yz$, $b=y^2+xz$, and
\[
\bar S=S/(ab),\quad \bar S_a=S/(a),\quad \bar S_b=S/(b),\quad \bar S_{ab}=S/(a,b).
\]
The two conics give an exact sequence
\[
0\longrightarrow\bar S\longrightarrow\bar S_a\oplus\bar S_b
\longrightarrow\bar S_{ab}\longrightarrow0.
\]
On the first conic, use the parametrization
$(y,x,z)=(\mathsf u^2,\mathsf u\mathsf v,\mathsf v^2)$.
The parameters $\mathsf u,\mathsf v$ are unrelated to $u,w$ in $R$.
A monomial surviving $(x^q,y^q,z^q)$ has both parameter exponents at
most $4Q-1$, and at least one at most $2Q-1$.
Its parameter degree is therefore at most $6Q-2$.
Thus $\bar S_a/(x^q,y^q,z^q)$ vanishes in all ordinary degrees at least
$3Q$. The same holds for $\bar S_b$.

Apply the Koszul complex on $\mathbf x^q=x^q,y^q,z^q$ to the exact
sequence. In degree $j$ there are no first boundaries, since
$j-2q=1-Q<0$. The scheme $\Proj\bar S_{ab}$ consists of four reduced
points: $P_0=(0,0,1)$ and $(\omega,\omega^2,1)$ with $\omega^3=1$.
Evaluation identifies $(\bar S_{ab})_n$ with $k^4$ for $n\geq2$.
Let $X,Y$ be the evaluations of $x,y$, let $\lambda$ sum the four
coordinates, and put $V_0=\ker\lambda$. The space $V_0$ is the
image of the linear forms. We have
\[
H_1(\mathbf x^q,\bar S_{ab})_j
=\{(a,b,c)\in(k^4)^3:X^qa+Y^qb+c=0\}\simeq(k^4)^2.
\]
The syzygies $(y^Q,z^Q,0)$ and $(z^Q,0,y^Q)$ on the first conic,
multiplied by linear forms, give the six-dimensional subspace
\[
(a,b)=(Y^QL_1+L_2,L_1),\qquad L_1,L_2\in V_0.
\]
The quotient is detected by $\lambda(b)$ and $\lambda(a+Y^Qb)$.

On the second conic, the syzygy $(z^Q,x^Q,0)$ multiplied by a linear
form $L$ gives $(a,b)=(L,X^QL)$.
Its quotient coordinates are $\lambda(X^QL)$ and $L(P_0)$, because
$XY$ is zero at $P_0$ and one at the other three points.
Taking $L=z$ gives $(0,1)$.
Taking $L=y$ when $Q\equiv1\pmod3$, or $L=x$ when
$Q\equiv2\pmod3$, gives $(1,0)$.
Thus the two conic images span $H_1(\mathbf x^q,\bar S_{ab})_j$.
The homology sequence now gives
$(\bar S/(x^q,y^q,z^q))_j=0$, proving surjectivity.

To count the character dimensions, the formal character series of
$A_q$ is $(1-\mathsf t^{3q})/(1-\mathsf t^3)$.
Here complex cube roots are used only to count monomial weights.
Its coefficients in degrees $3Q-3$ and $3Q+1$ are one and zero,
whereas the ordinary dimensions are $3Q^2-2$ and $3Q^2-6$.
Interchanging $x,y$ equates the two nonzero character dimensions.
They are therefore $Q^2-1$ and $Q^2-2$, as claimed.
\end{proof}

\subsection{The kernel line bundle and local cohomology}
\begin{lemma}\label{lem:cofactor}
Let $S'=k[u,w]$ and let $M$ be an $n\times(n+1)$ matrix of linear forms
of generic rank $n$. If $G$ is the gcd of its nonzero maximal minors, then
$\ker_{S'} M=S'p$ for a primitive homogeneous vector $p$, and
$\deg p=n-\deg G$.
\end{lemma}
\begin{proof}
The cofactor vector $\mathbf c$ satisfies $M\mathbf c=0$. Dividing by $G$ gives
$p$ with relatively prime coordinates. A polynomial kernel vector is $ap$
for some $a\in k(u,w)$. For each irreducible $\pi\in S'$, one coordinate
of $p$ has $\pi$-valuation zero, so $a$ has nonnegative $\pi$-valuation.
Since $S'$ is a UFD, $a\in S'$.
\end{proof}
Apply Lemma~\ref{lem:cofactor} to $M_q$. Let $G_q$ be the homogeneous
gcd of its nonzero maximal minors, and let $c_q$ be the coefficient
degree of its primitive kernel vector. Thus
\begin{equation}\label{eq:index-gcd}
c_q=n_q-\deg G_q.
\end{equation}
Sheafification on the parameter line gives an exact sequence
\begin{equation}\label{eq:kernel-sheaf}
0\longrightarrow\mathcal K_q
\longrightarrow\mathcal O_{\mathbb P^1}(-1)^{n_q+1}
\longrightarrow\mathcal O_{\mathbb P^1}^{n_q}
\longrightarrow\mathcal T_q\longrightarrow0,
\end{equation}
where
\[
\mathcal K_q\simeq\mathcal O_{\mathbb P^1}(-c_q-1),
\qquad \length\mathcal T_q=\deg G_q=n_q-c_q.
\]
To see the length assertion, work over an algebraic closure.
At a point $a\in\mathbb P^1$, Smith normal form over the local discrete
valuation ring gives
\[
(\mathcal T_q)_a\simeq
\bigoplus_j\mathcal O_{\mathbb P^1,a}/(t^{a_j}),\qquad a_j>0.
\]
Its length is $\sum_j a_j$, the local order of $G_q$.


The torsion sheaf $\mathcal T_q$ records rank drops on the parameter
line. It is different from the finite-length graded module
$H^0_{(u,w)}(\mathcal N_q)$, which sheafifies to zero.
The next lemma recovers that module from the kernel degree.

\begin{lemma}\label{lem:duality}
Let $S'=k[u,w]$, let $\mathfrak n=(u,w)$, and let
$M:S'(-1)^{n+1}\to S'^{\,n}$ have generic rank $n$.
Suppose its primitive kernel vector $p=(p_1,\ldots,p_{n+1})$ has
coefficient degree $c>0$, and put $N=\coker M$.
For a graded $k$-vector space $V$, write
$(V^\vee)_d=\Hom_k(V_{-d},k)$. Then
\begin{equation}\label{eq:pencil-duality}
H^0_{\mathfrak n}(N)^\vee
\simeq \bigl(S'/\mathfrak n^c\bigr)(c-1).
\end{equation}
In particular,
\[
\dim_k H^0_{\mathfrak n}(N)_d=
\begin{cases}c-d,&0\leq d\leq c-1,\\0,&\text{otherwise},\end{cases}
\qquad
\dim_k H^0_{\mathfrak n}(N)=\frac{c(c+1)}2.
\]
Its annihilator is $\mathfrak n^c$. Thus its Loewy length is $c$, and
$w^{c-1}H^0_{\mathfrak n}(N)\ne0$.
\end{lemma}
\begin{proof}
Lemma~\ref{lem:cofactor} gives the resolution
\[
0\longrightarrow S'(-c-1)\xrightarrow{p}S'(-1)^{n+1}
\xrightarrow{M}S'^{\,n}\longrightarrow N\longrightarrow0.
\]
Let $L=(p_1,\ldots,p_{n+1})$. Applying
$\Hom_{S'}(-,S'(-2))$ and graded local duality gives
\[
H^0_{\mathfrak n}(N)^\vee\simeq(S'/L)(c-1).
\]
All generators of $L$ have degree $c$, so $L_d=0$ for $d<c$.
Also $N_d=0$ for $d<0$, and the same holds for its local-cohomology
submodule. The displayed duality therefore forces
$(S'/L)_a=0$ for every $a\geq c$. Hence $L=\mathfrak n^c$.
This proves~\eqref{eq:pencil-duality} and the Hilbert function.
The graded dual of a finite-dimensional module has the same annihilator.
Thus $\Ann_{S'}H^0_{\mathfrak n}(N)=\mathfrak n^c$, and the remaining
claims follow.
\end{proof}

\begin{proposition}\label{prop:regularity-reduction}
If $q\geq8$ and $c_q>0$, then $\reg_R I^{[q]}\geq3q/2+c_q+1$.
\end{proposition}
\begin{proof}
Put $T_q=R/I^{[q]}=A_q[u,w]/(uz^4+wh)$, and give $x,y,z$ degree $(1,0)$
and $u,w$ degree $(0,1)$. The part of $T_q$ of first degree $j=3q/2+1$ and character one
is $\mathcal N_q$, with its second grading over $S'=k[u,w]$.
Lemma~\ref{lem:duality} gives a nonzero class of second degree $c_q-1$,
hence total degree $3q/2+c_q$. Since $x,y,z$ act nilpotently on $T_q$,
this class belongs to $H^0_\mm(T_q)$. Cohen--Macaulayness of $R$ gives
$H^1_\mm(I^{[q]})\simeq H^0_\mm(T_q)$, proving the bound.
\end{proof}

For the matrix calculations below, $\Id$ denotes an identity matrix.
We use $\mathcal V$ for a discrete valuation ring $k\llbracket t\rrbracket$
or $k\llbracket s\rrbracket$. The notation $[t^N]g$ means coefficient
extraction, and $[m]g$ denotes the coefficient of a monomial $m$.

\section{Escape parameters and critical syzygies}\label{sec:escape}
We relate the central pencil to syzygy bundles on the quartics
$C_\alpha=V(h+\alpha z^4)$. 

\subsection{Syzygies on the smooth quartics}\label{subsec:quartic-syzygies}
Syzygy bundles encode relations among the generators on a projective
curve. Their stability and its behaviour under Frobenius enter the
study of Hilbert--Kunz multiplicity \cite{BrennerStable,Trivedi}.
Here we need the splitting of a specific Frobenius pullback and the
order to which its sections lift along the quartic family.
Put $S=k[x,y,z]$, $A=x^2+yz$, and $B=y^2+xz$, so $h=AB$.
Let $\alpha\ne0$ and $C=C_\alpha$. On the chart $z=1$, the partial
derivatives of $h+\alpha z^4$ in $x,y$ are $A,B$.
They cannot both vanish because $AB=\alpha\ne0$.
The only points at infinity are $P_x=(1,0,0)$ and $P_y=(0,1,0)$.
The $z$ derivative $x^3+y^3$ is nonzero at both. Thus $C$ is smooth,
and adjunction and the intersection with $z=0$ give
\begin{equation}\label{eq:quartic-canonical}
\omega_C\simeq\mathcal O_C(1),\qquad
\mathcal O_C(1)\simeq\mathcal O_C(2P_x+2P_y).
\end{equation}
In particular, $C$ has genus three.

For dyadic $q\geq2$, set
\[
E_q=\Syz_C(x^q,y^q,z^q)(3q/2),
\]
where the untwisted syzygy bundle is the kernel of the surjection
$\mathcal O_C(-q)^3\to\mathcal O_C$.
The bundle $E_q$ has rank two and trivial determinant.
For $m=2^a$, one has $E_{qm}\simeq(F_C^a)^*E_q$.
This relation will turn one critical splitting into splittings at all
higher dyadic multiples.

On $z=1$, the functions $A,B$ are units with $AB=\alpha$.
All Laurent expansions are in $A$, with $A^{-1}=\alpha^{-1}B$.
The following correspondence connects sections of $E_q$ with the
kernel and cokernel of the multiplication pencil.
\begin{lemma}\label{lem:ann-syz}
If $f\in S_4$ has character zero, $C=V(f)$ is smooth, and $N<2q$, then
\[
(0:_{A_q}f)_{N-4,i}\simeq
H^0(C,\Syz_C(x^q,y^q,z^q)(N))_i.
\]
For $q\geq4$ and $f=h+\alpha z^4$, the central kernel is $H^0(C,E_q(1))_1$ and
the dual central cokernel is $H^0(C,E_q)_2$.
\end{lemma}
\begin{proof}
The restriction sequence
\[
0\longrightarrow\mathcal O_{\mathbb P^2}(j-4)
\xrightarrow{\,f\,}\mathcal O_{\mathbb P^2}(j)
\longrightarrow\mathcal O_C(j)\longrightarrow0
\]
and $H^1(\mathbb P^2,\mathcal O(j-4))=0$ show that each coefficient
section of a syzygy lifts to a polynomial. Thus
\[
x^qa+y^qb+z^qc=f\widetilde F,
\qquad \deg\widetilde F=N-4.
\]
Changing a coefficient lift by a multiple of $f$ changes
$\widetilde F$ by an element of $(x^q,y^q,z^q)$.
Conversely, an annihilator class gives such a relation and hence a
syzygy on $C$.

If the annihilator class is zero, subtract multiples of $f$ from the
coefficient lifts to obtain a polynomial syzygy on $x^q,y^q,z^q$.
This regular sequence has no nonzero syzygy of total degree below $2q$.
The correspondence is therefore an isomorphism. It respects characters
because $f$ has character zero. The central cokernel statement follows
from the perfect socle pairing on $A_q$, whose socle has degree $3q-3$
and character zero.
\end{proof}
\subsection{Escape polynomials}
For the line-$S_4$ quartics, the escape map is
$s\mapsto s^4+ts^{-4}$, starting at $1$.
See~\cite[Definitions 2.1--2.2]{Borevitz} and~\cite[Theorem 5.10]{Monsky}.
Writing its iterates as rational functions gives the polynomials

\[
P_1=1+t,\quad Q_1=1,\qquad
P_{j+1}=P_j^8+tQ_j^8,\quad Q_{j+1}=P_j^4Q_j^4.
\]
\begin{proposition}\label{prop:escape}
The $P_j$ are squarefree and pairwise coprime, with
$\deg P_j=8^{j-1}$ and $\deg Q_j=8^{j-1}-4^{j-1}$.
Their roots are nonzero and give the first zero of the iteration
$s\mapsto s^4+ts^{-4}$ starting at $1$.
\end{proposition}
\begin{proof}
Induction gives $P_j(0)=Q_j(0)=1$ and $\gcd(P_j,Q_j)=1$.
At a root of $P_j$, the next numerator is nonzero and the next
denominator is zero. All later numerators remain nonzero there.
Hence the $P_j$ are pairwise coprime.
The identity $P'_{j+1}=Q_j^8$ proves squarefreeness.
The degree formulas follow from the recurrences, and
$P_j/Q_j$ are the stated iterates.
\end{proof}
\subsection{Critical syzygies}
Let $\alpha$ be a root of $P_e$, put $q_0=2\cdot4^e$, $\Phi=q_0/4$,
$n=2\Phi-1$, $A=x^2+yz$, $B=y^2+xz$, and $f=AB+\alpha z^4$.
We work over an algebraic closure when computing local exponents.

We will use the translations
\[
g_\sigma(x,y,z)=(x+\sigma z,y+\sigma^2z,z),
\qquad \sigma\in\mathbb F_4.
\]
They preserve $A$, $B$, and $f$. Their induced action on a syzygy triple is
\[
(a,b,c)\longmapsto
\bigl(g_\sigma a,g_\sigma b,
g_\sigma c+\sigma^q g_\sigma a+\sigma^{2q}g_\sigma b\bigr).
\]
Indeed this formula follows by applying $g_\sigma$ to
$x^qa+y^qb+z^qc=0$.
For a fixed triple, $a,b$ are invariant. If $q=4^j$, then
$c+ax+by$ is invariant. If $q=2\cdot4^j$, then $c+ay+bx$ is invariant.
\begin{lemma}\label{lem:pade}
Put $\kappa_e=(\Phi+1)/3$. Let $\Sigma$ be the upper shift matrix of size $\kappa_e$,
$\mathsf J$ the basis reversal, $\Id$ the identity, and
\[
\mathcal H_e(\alpha,b)=\Id+\sum_{i=0}^{e-2}\Sigma^{(\Phi-2\cdot4^i)/3}
+b\mathsf J\diag(\alpha,\alpha^4,\ldots,\alpha^{\Phi-1}).
\]
If $P_e(\alpha)=0$, then $\mathcal H_e(\alpha,1)$ has a
one-dimensional kernel, and its determinant has a simple zero at
this parameter. A nonzero kernel vector
$\varphi=(\varphi_{1+3j})_{0\leq j<\kappa_e}$ has both endpoint
coordinates nonzero and satisfies $\varphi_{1+3j}=0$ for odd $j$.
Thus only odd exponents occur in
$\varphi(A)=\sum_a\varphi_aA^a$.
\end{lemma}
\begin{proof}
For $e=1$, the matrix is $[1+b\alpha]$, and the assertions follow
at the root $\alpha=1$ of $P_1$.
For $e\geq2$ set $\Phi'=\Phi/4$, $r=(\Phi'+1)/3$, so $\kappa_e=4r-1$ and $\Phi'=3r-1$.
Partition the indices into consecutive groups of sizes $r,2r-1,r$. All nonidentity
shifts are at least $\Phi'$, so the only off-diagonal shift block is the
first-to-last block $\mathsf T=\mathcal H_{e-1}(\alpha,0)$. If $\mathsf B'=\mathsf
J\diag(\alpha,\alpha^4,\ldots,\alpha^{\Phi'-1})$, then $(\mathsf
B')^2=\alpha^{\Phi'}\Id$. The outer matrix is
\[
\begin{pmatrix}\Id&\mathsf T+b\alpha^{3\Phi'}\mathsf B'\\b\mathsf B'&\Id\end{pmatrix},
\]
and the middle determinant is $(1+b\alpha^{\Phi/2})^{2r-1}$. A Schur complement gives
\[
\det \mathcal H_e(\alpha,b)=
(1+b\alpha^{\Phi/2})^{2r-1}b^r\alpha^{\Phi'r/2}
\det \mathcal H_{e-1}\left(\alpha,\frac{1+b^2\alpha^\Phi}{b\alpha^{\Phi'}}\right).
\tag{P}
\]
The identity (P) holds with $\alpha,b$ as indeterminates.
To evaluate its order at a root of $P_e$, put
\[
\tau_0=\alpha,\qquad
\tau_{i+1}=\alpha^3(\tau_i^4+\tau_i^{-4}).
\]
At stage $j=e-i$, the second argument is
$(\tau_i/\alpha)^{4^{j-1}}$. Substitution in (P) verifies this formula.
For $i\geq1$, the escape recurrence gives
\[
\tau_i=\frac{P_i^8}{\alpha Q_i^8},\qquad
1+\tau_i=\frac{P_{i+1}}{\alpha Q_i^8},
\]
and $1+\tau_0=P_1$.
Pairwise coprimality shows that all nonterminal factors are units
at a root of $P_e$. The terminal factor $1+\tau_{e-1}$ has order one
because $P_e$ is squarefree. Hence the determinant has a simple zero.
Corank at least two would force every cofactor, and therefore the
determinant derivative, to vanish. The corank is exactly one.

We also track the kernel through these invertible block steps. Its middle coordinates
vanish. If $\xi,\zeta$ are its first and last $r$ coordinates, then $\zeta=b\mathsf
B'\xi$. The vector $\eta=\mathsf B'\xi$ lies in the smaller kernel with the
transformed parameter in (P), and $\xi=\alpha^{-\Phi'}\mathsf B'\eta$, $\zeta=b\eta$.
Consequently
\[
\xi_0=\alpha^{-1}\eta_{r-1},\qquad\zeta_{r-1}=b\eta_{r-1}.
\]
Starting with the nonzero one-dimensional base vector, induction proves the two
endpoint assertions at every intermediate stage. The integer $r$ is odd, so reversal
preserves index parity, and the last block starts at the even index $\Phi'$.
Induction also shows that only even indices survive, yielding the asserted odd
exponents.
\end{proof}

\begin{theorem}\label{thm:critical}
With the natural twists,
\[
E_{q_0}:=\Syz_C(x^{q_0},y^{q_0},z^{q_0})(3q_0/2)
\simeq\mathcal O_C(1)\oplus\mathcal O_C(-1).
\]
The critical annihilator in degree $3q_0/2-5$ is one-dimensional and has auxiliary
character zero. At every dyadic $m\geq1$, the central weight-one pencil at $q=q_0m$
has exactly $\lfloor4m/3\rfloor$ positive local elementary divisors at this
parameter.
\end{theorem}
\begin{proof}
\emph{Step 1: construction of the critical section.}
We use the two points at infinity and the line-bundle identities
in~\eqref{eq:quartic-canonical}.

On $z=1$, one has $B=\alpha/A$ and $x^4+x=A^2+B$.
The affine ring is free on $1,x,y,xy$ over $k[A,A^{-1}]$:
use $y=A+x^2$ and the monic degree-four equation for $x$.
The pole orders at $(P_x,P_y)$ are
\[
A:(4,-4),\qquad x:(2,1),\qquad y:(1,2),\qquad xy:(3,3).
\]
In the chart $x=1$, the parameter $s=y/x$ gives $z=s^2$ times a unit.
Interchanging $x,y$ gives the orders at $P_y$.

Repeated fourth powers give
\[
x^{q_0}=y+U_e,\quad y^{q_0}=x+V_e,\quad
U_e=\sum_{i=0}^{e}A^{4^i}+\sum_{i=0}^{e-1}B^{2\cdot4^i},\quad V_e=U_e(B,A).
\]
Let $\varphi$ be the polynomial from Lemma~\ref{lem:pade}, and put
$\psi=\varphi(B)U_e+\varphi(A)V_e$.
This Laurent polynomial is invariant under $A\leftrightarrow B$.
Its exponents are divisible by three and lie between
$-2\Phi+1$ and $2\Phi-1$.
For $a=1,4,\ldots,\Phi-1$, its coefficient at $A^{\Phi+a}$ is
\[
\varphi_a+\sum_{\substack{0\leq i<e-1\\a<2\cdot4^i}}\varphi_{\Phi+a-2\cdot4^i}+\alpha^{\Phi-a}\varphi_{\Phi-a},
\]
where an absent coefficient of $\varphi$ is zero.
The first two terms come from the positive powers $A^\Phi$ and
$A^{2\cdot4^i}$ in $V_e$. The last comes from
$A^{2\Phi}B^{\Phi-a}=\alpha^{\Phi-a}A^{\Phi+a}$ in
$\varphi(B)U_e$. All other terms have exponent less than $\Phi$.
Thus these coefficients form the vector
$\mathcal H_e(\alpha,1)\varphi$.

For example, when $e=2$, one has $\Phi=8$ and
\[
\mathcal H_2(\alpha,1)=
\begin{pmatrix}
1&0&1+\alpha^7\\
0&1+\alpha^4&0\\
\alpha&0&1
\end{pmatrix}.
\]
Its three rows are the coefficients of $A^9,A^{12},A^{15}$ in
$\psi$, in this order. The same exponent comparison gives the formula
for every $e$.
These coefficients vanish by Lemma~\ref{lem:pade}. Reciprocity gives
vanishing of the high negative coefficients. Therefore
$\supp\psi\subseteq[-\Phi+1,\Phi-1]$.

The triple
\[
a=\varphi(B),\qquad b=\varphi(A),\qquad
c=\varphi(A)x+\varphi(B)y+\psi
\]
satisfies $x^{q_0}a+y^{q_0}b+c=0$.
Its poles are bounded by $2n$ at both points at infinity.
Since $\mathcal O_C(n)=\mathcal O_C(2nP_x+2nP_y)$, it homogenizes
to a syzygy whose coefficients have degree $n$.

\emph{Step 2: the section has no zeros.}
Suppose $\varphi(A)$ and $\varphi(\alpha/A)$ have a common nonzero root.
Let
\[
g(A)=\gcd(\varphi(A),A^{\Phi-1}\varphi(\alpha/A)),\qquad \rho=\deg g>0,
\]
with $g$ monic. The second polynomial has nonzero constant term
$\varphi_{\Phi-1}\alpha^{\Phi-1}$, so $g(0)\ne0$.
The roots of $g$, with multiplicities, are stable under
$A\mapsto\alpha/A$ and multiplication by cube roots of unity.
Every nonzero root has even multiplicity because $\varphi(A)/A$
is a polynomial in $A^2$.
Each cube-root orbit has length three. Hence $6\mid\rho$.

Put $\Gamma(A)=A^{-\rho/2}g(A)$. Its character is zero: $g$ has exponents divisible
by three and $\rho/2$ is divisible by three. Reciprocity gives
$\Gamma(\alpha/A)=\sigma\Gamma(A)$ for a nonzero scalar $\sigma$. Applying the
involution twice gives $\sigma^2=1$, hence $\sigma=1$ in characteristic two.
Therefore
\[
\varphi'(A)=\varphi(A)/\Gamma(A)=A^{\rho/2}\varphi(A)/g(A)
\]
is a nonzero polynomial with positive exponents congruent to one modulo three, but
degree $\Phi-1-\rho/2<\Phi-1$. Because $\Gamma$ is invariant under reciprocity,
division of $\psi$ by $\Gamma$ is exact in the Laurent ring and gives
$\varphi'(B)U_e+\varphi'(A)V_e$. Its support is contained in
$[-\Phi+1+\rho/2,\Phi-1-\rho/2]$: leading and trailing Laurent exponents add under
multiplication, and $\Gamma$ has extremal exponents $\pm \rho/2$. Thus $\varphi'$
defines another kernel vector of the same matrix $\mathcal H_e$, with zero last
coordinate. This contradicts one-dimensionality and the nonzero last coordinate of
the original vector. There is no affine base point.

At infinity, the homogenizations of $a,b$ are divisible by $z$.
The terms of degree $n$ in the third coefficient are
$\varphi_{\Phi-1}(A^{\Phi-1}x+B^{\Phi-1}y)$.
All its other terms vanish at $P_x,P_y$, and its value at each point
is $\varphi_{\Phi-1}\ne0$. Thus the section is nowhere zero.

\emph{Step 3: splitting and Frobenius pullback.}
The constructed section is a nowhere-zero section of $E_{q_0}(-1)$. It gives an exact sequence
\[
0\longrightarrow\mathcal O_C(1)\longrightarrow E_{q_0}\longrightarrow\mathcal
O_C(-1)\longrightarrow0.
\]
Serre duality gives $H^1(C,\mathcal O_C(2))=0$, so the sequence splits.
By Lemma~\ref{lem:ann-syz}, the critical annihilator is
\[
H^0(C,E_{q_0}(-1))
=H^0(C,\mathcal O_C\oplus\mathcal O_C(-2))=k.
\]
The constructed relation has total character zero, so its generator
$F_0$ has character zero. Each translation and the interchange of
$x,y$ act trivially on this line: their actions have order two,
and the only scalar with square one in characteristic two is one.

Absolute Frobenius pullback gives $E_{q_0m}\simeq\mathcal O_C(m)\oplus\mathcal
O_C(-m)$. The positive summand is generated by the $m$th power of the character-zero
critical section. Writing $\widetilde u_q=f^{m-1}F_0^m$, the annihilator--syzygy
correspondence gives central primal kernel and dual obstruction space
\[
\widetilde u_q(S/(f))_{m+1,1},\qquad
\widetilde u_q(S/(f))_{m,2}.
\]
For $m\geq2$ the other summand contributes no sections to the first space. For $m=1$
its possible constant section has character zero and contributes none to character
one. This last character assertion also follows from the natural determinant
character, which is zero, and the character-zero positive line. The negative summand
contributes no sections to the dual space for any $m\geq1$.

Finally $\dim_k(S/(f))_{m,2}=\lfloor4m/3\rfloor$ for dyadic $m$. One can check this
from the ordinary Hilbert dimension $4m-2$ for $m\geq2$ and the formal character
series $(1-\mathsf t^4)/(1-\mathsf t^3)$: its coefficient is the difference of the
weight-zero and weight-two dimensions, while the weight-one and weight-two dimensions
coincide by symmetry. The cases $m=1,2$ are direct. This proves the stated number of
positive local factors.

\end{proof}

\section{Finite determinantal support}\label{sec:support}
Write $q=2^n\geq4$. For $\alpha\ne0$ let $\ell(\alpha)$ be the first
$j$ with $P_j(\alpha)=0$, with value $\infty$ if there is no such $j$.
Monsky's Hilbert--Kunz formula for the line-$S_4$ quartics
\cite[Corollary 3.15 and Theorem 5.8]{Monsky}, in the escape notation of \cite{Borevitz}, gives
\begin{equation}\label{eq:HK-fibre}
\dim_{\bar k}\bigl(A_q/(h+\alpha z^4)\bigr)=
\begin{cases}
3q^2-4,&n\leq2\ell(\alpha),\\
3q^2+4^{n-2\ell(\alpha)},&n>2\ell(\alpha).
\end{cases}
\end{equation}
Here $A_q$ denotes its scalar extension to $\bar k$. The generic parameter never escapes
because every $P_j$ is a nonzero polynomial over $\mathbb F_2$.

\begin{theorem}\label{thm:support}
For dyadic $q\geq4$, the central weight-one cokernel at a finite nonzero parameter
$\alpha$ is nonzero precisely when
\[
P_j(\alpha)=0\quad\text{for some }j\geq1\text{ with }q\geq2\cdot4^j.
\]
The parameter zero does not belong to the support.
\end{theorem}
\begin{proof}
For each degree and character, the fibre of $A_q/(h+\alpha z^4)$
is the cokernel of a matrix with polynomial entries in $\alpha$.
Its dimension is at least the corresponding generic dimension.
If $n\leq2\ell(\alpha)$, the sum of these dimensions equals the generic
sum by \eqref{eq:HK-fibre}. Equality therefore holds in every summand.
In particular the central weight-one cokernel vanishes, since its
generic dimension is zero by Proposition~\ref{prop:rank}.

If $j=\ell(\alpha)<\infty$ and $n>2j$, write $q=(2\cdot4^j)m$.
Theorem~\ref{thm:critical} gives central cokernel dimension
$\lfloor4m/3\rfloor>0$. At $\alpha=0$, Proposition~\ref{prop:rank}
gives surjectivity directly.
\end{proof}

\section{The first obstruction}\label{sec:obstruction}
Knowing the support of the torsion sheaf does not determine its length.
We compute the local exponents by lifting a critical syzygy along the
one-parameter deformation $f_t=f_\alpha+tz^4$.
A lift modulo $t^N$ is a syzygy over the infinitesimal parameter
neighbourhood $\Spec k[t]/(t^N)$. Its first obstruction detects a local
Smith exponent.

Put $S=k[x,y,z]$, $f=f_\alpha=h+\alpha z^4$, and $f_t=f+tz^4$.
Write $\Lambda_q$ for extraction of the coefficient of $(xyz)^{q-1}$ in $A_q$.
Keep $q=q_0=2\cdot4^e$, $d_0=3q/2-5$, and a critical generator $F_0\in(A_q)_{d_0}$.
Theorem~\ref{thm:critical} proves that this annihilator is one-dimensional, character
zero, and translation invariant.

\begin{lemma}
If a class $\Xi\in(A_q)_{3q-10}$ has character zero and is translation invariant, then
\[
p\longmapsto\Lambda_q(z^4\Xi p)\quad(p\in S_3)
\]
is a scalar multiple of $\theta_0=[x^3]+[y^3]+[xyz]$. If $\Xi$ is a square, that
scalar is zero.
\end{lemma}
\begin{proof}
For an invariant $\Psi$ of degree $3q-5$, invariance of $\Lambda_q(\Psi xy)$ under
$g_\sigma$ gives
\[
\sigma^2\Lambda_q(\Psi xz)+\sigma\Lambda_q(\Psi yz)+\sigma^3\Lambda_q(\Psi z^2)=0\quad(\sigma\in\mathbb F_4).
\]
A polynomial of degree at most three vanishing at all four elements is zero, so all
three coefficients vanish. Apply this to $\Psi=z^3A\Xi,z^3B\Xi,z^5\Xi$. It gives
$\theta(Ax)=\theta(By)=\theta(z^3)=0$. Character zero permits nonzero values only on
$x^3,y^3,xyz,z^3$, proving the formula. For a square $\Xi$, its value at $x^3$ is
zero by parity of the $y$ exponent, since $q-1$ is odd.
\end{proof}

\begin{proposition}\label{prop:first-lift}
There is a character-zero, translation-invariant $F_1$ of degree $d_0$ with
$fF_1=z^4F_0$ in $A_q$. For such a lift,
$\Lambda_q(z^4F_0F_1p)=\gamma\theta_0(p)$, and the scalar is independent of its choice.
\end{proposition}
\begin{proof}
The nowhere-zero critical section gives an equivariant sequence
\[0\longrightarrow\mathcal O_C\longrightarrow\Syz_C(x^q,y^q,z^q)(3q/2-1)\longrightarrow\mathcal O_C(-2)\longrightarrow0.\]
The determinant has its natural action: the coordinate representation
has determinant one, and the critical section is invariant.
After twisting by three, the subbundle and quotient are
$\mathcal O_C(3)$ and $\mathcal O_C(1)$.
Their section spaces are $S_3$ and $S_1$, and
$H^1(C,\mathcal O_C(3))=0$.
Lemma~\ref{lem:ann-syz} therefore gives the equivariant exact sequence
\[
0\to F_0S_3\to \mathcal K\to S_1\to0,\qquad \mathcal K=\Ann_{A_q}(f)_{d_0+3}.
\]
By the perfect socle pairing, existence of $F_1$ is equivalent to vanishing on
$\mathcal K$ of $v\mapsto\Lambda_q(z^4F_0v)$. The invariant-cubic lemma applied to
$F_0^2$ makes it zero on $F_0S_3$. It factors to an invariant functional on $S_1$,
which must vanish: the cube-root action eliminates its values at $x,y$, and
$g_1(x)=x+z$ then eliminates its value at $z$.

Project any lift onto character zero. Its translation defect lies in $kF_0$, say
$g_\sigma F_1-F_1=\lambda_\sigma F_0$. The scalars are additive in $\sigma$ and are
equal for the three nonzero $\sigma$, by conjugation with the cube-root action.
Additivity forces all to be zero. Thus the lift is invariant. The cubic lemma applies
to $F_0F_1$. Two lifts differ by a multiple of $F_0$ and $z^4F_0^2=(fF_1)F_0=0$,
proving independence.
\end{proof}

\begin{theorem}\label{thm:critical-order}
At a root of $P_e$ the critical multiplication map has one positive local exponent,
and it is exactly two. Equivalently the scalar $\gamma$ in the preceding proposition
is nonzero. The central pencil at the critical level also has its sole positive
exponent equal to two.
\end{theorem}
\begin{proof}
Suppose a lift modulo $t^3$ exists. Choose character-zero coefficients $F_0,F_1,F_2$ with
$fF_i=z^4F_{i-1}$ and let $\beta=\alpha+t$ in $\Omega=k[t]/(t^3)$. The same
translation-defect argument makes each $F_i$ invariant, inductively. Set $\widetilde
F=F_0+tF_1+t^2F_2$ and choose representatives with each coordinate exponent below
$q$. There are unique coefficient polynomials of degree $n=q/2-1$ satisfying
\[
f_\beta \widetilde F=x^qa+y^qb+z^qc.
\]
The decomposition is unique even over $\Omega$: its total degree is
$3q/2-1<2q$, so no monomial is divisible by two of $x^q,y^q,z^q$.
For each translation, $g_\sigma\widetilde F-\widetilde F$ belongs
to $(x^q,y^q,z^q)$. Multiplying this difference by $f_\beta$ and
using uniqueness shows that the corresponding coefficient differences
are multiples of $f_\beta$.
Consequently the coefficient sections on the relative curve satisfy
\[
g_\sigma a=a,\qquad g_\sigma b=b,\qquad
g_\sigma c=c+\sigma^2a+\sigma b,
\]
because $\sigma^q=\sigma^2$ and $\sigma^{2q}=\sigma$.
It follows that $\psi=c+ay+bx$ is invariant.

The relative affine coordinate ring is free on $1,x,y,xy$ over $\Omega[A,A^{-1}]$,
$AB=\beta$. For $a_0+a_1x+a_2y+a_3xy$ in this basis, invariance under
$g_1$ forces $a_3=0$ and $a_1=a_2$. Invariance under $g_\omega$
then gives $(\omega+\omega^2)a_1=0$. Since $\omega+\omega^2=1$,
the invariant ring is $\Omega[A,A^{-1}]$. Consequently $a,b,\psi$ are
scalar Laurent polynomials.

Their supports lie in $[-\Phi+1,\Phi-1]$, $\Phi=q/4$. This pole assertion is valid
even with nilpotent coefficients: at $P_x$ use the parameter $s=y/x$ in
$s^2+(1+s^3)z+sz^2+\beta z^4=0$. The $z$ derivative is a unit there, so the completed
ring is $\Omega\llbracket s\rrbracket$ and $z=s^2$ times a unit. The top Laurent term
of $A$ has pole order four with unit coefficient. A nonzero top Laurent coefficient,
even if nilpotent, cannot cancel a lower pole order. The homogeneous degree bound
$n=2\Phi-1$ gives the asserted scalar support, and symmetry treats the other point.
The same bound holds for $\psi$ because $a,b$ have poles at most $4\Phi-4$, so
multiplying by $x,y$ adds at most two.

The exact power identities give $\psi=aU_e+bV_e$. The leading term of $U_e$ is
$A^{2\Phi}$ while that of $V_e$ is $A^\Phi$. Any nonnegative exponent in $a$ would
produce an uncancellable term of degree at least $2\Phi$, because $b$ has maximum
exponent $\Phi-1$. The reciprocal argument excludes nonpositive exponents in $b$.
Their characters force
\[
a=\sum_{j=0}^{\kappa_e-1}p_jB^{1+3j},\qquad
b=\sum_{j=0}^{\kappa_e-1}r_jA^{1+3j}.
\]
Let $\mathcal H_e(\beta,1)=\mathsf L_e+\mathsf B_e$ with $\mathsf L_e$ its upper
unipotent part. The high positive and negative equations are exactly
\[
\mathsf L_er+\mathsf B_ep=0,\qquad \mathsf B_er+\mathsf L_ep=0.
\]
Thus $(r,p)$ is a kernel vector modulo $t^3$ of the square matrix
$\mathcal M_e=\left(\begin{smallmatrix}\mathsf L_e&\mathsf B_e\\\mathsf B_e&\mathsf L_e\end{smallmatrix}\right)$.
Constant block operations give $\det\mathcal M_e=(\det \mathcal H_e)^2$, which has
order two at the chosen root by Lemma~\ref{lem:pade}. Every Smith exponent of this
square matrix is at most two. Therefore every kernel vector modulo $t^3$ reduces to
zero modulo $t$. But $(r,p)$ does not: otherwise $a,b,\psi,c$ would all reduce to
zero as sections on $C$, and the annihilator--syzygy correspondence would give
$F_0=0$. This contradiction rules out a second lift. The first lift exists, so the
critical exponent is exactly two.

If $\gamma=0$, the same annihilator pairing used to construct $F_1$, now applied to
$z^4F_1$, would give $F_2$: the invariant functional vanishes on $F_0S_3$ and then on
the quotient $S_1$. Hence $\gamma\ne0$. Finally the order-two central residual
pairing is nonzero since $\theta_0(y\cdot xz)=1$. The critical geometric count gives
only one positive central factor, proving its exponent is two as well.
\end{proof}

\subsection{Residual pairings and coefficient extraction}
For a power of two $r$, let $\mathcal C_r$ send $x^iy^jz^k$ to
$x^{(i-r+1)/r}y^{(j-r+1)/r}z^{(k-r+1)/r}$ when all three exponents are congruent to
$r-1$ modulo $r$, and to zero otherwise. Extend it $k$-linearly: it leaves
coefficients unchanged. For a polynomial $p$ let $p^{\langle r\rangle}$ denote the
polynomial with the same monomials and coefficients raised to the $r$th power.

\begin{lemma}
For every dyadic $r$,
\[
\mathcal C_r(yf^{r-1})=\begin{cases}y&r\equiv1\pmod3,\\x&r\equiv2\pmod3,\end{cases}
\qquad \mathcal C_r(zf^{r-1})=z,
\]
and the analogous $x$ identity exchanges $x,y$. If $q=rq'$ and $\Lambda_n$ extracts
the coefficient of $(xyz)^{n-1}$, then for arbitrary polynomials $\Psi,p$,
\[
\Lambda_{rq'}(\Psi p^r)=\Lambda_{q'}\bigl(\mathcal C_r(\Psi)p^{\langle r\rangle}\bigr).
\tag{C}
\]
\end{lemma}
\begin{proof}
The identity (C) is a direct monomial comparison: a monomial of $p^r$ has exponent
vector $r\nu$, so its partner in $\Psi$ must have exponent vector
$r(q'-1-\nu)+(r-1,r-1,r-1)$. Its coefficient in $p^r$ is the $r$th power of the
original coefficient, accounting for $p^{\langle r\rangle}$.

Direct inspection of the five terms of $f$ gives $\mathcal C_2(xf)=y$, $\mathcal
C_2(yf)=x$, and $\mathcal C_2(zf)=z$, independently of $\alpha$. Also
\[
\mathcal C_2(\Psi G^2)=\mathcal C_2(\Psi)G^{\langle2\rangle},\qquad \mathcal C_{2s}=\mathcal C_s\circ \mathcal C_2.
\]
Since $f^{\langle2\rangle}=f_{\alpha^2}$, these identities yield
\[
\mathcal C_{2s}(yf_\alpha^{2s-1})=\mathcal C_s(xf_{\alpha^2}^{s-1}).
\]
Induction alternates $x$ and $y$ at each doubling. Powers of two alternate between
residues one and two modulo three. The $z$ identity follows by the same induction.
This proves the result over arbitrary fields of characteristic two. No inverse
Frobenius on coefficients is used.
\end{proof}

\begin{lemma}
Let $\theta_0=[x^3]+[y^3]+[xyz]$ on cubic forms and, for dyadic $m\geq2$, put
$\eta_m(\Psi)=\theta_0(\mathcal C_m(f^{m-1}\Psi))$ for forms $\Psi$ of degree $2m+1$. Then
\[
\Ann(\eta_m)=(A^m,B^m,z^{2m}):f^{m-1}
=(A^{m/2},B^{m/2},f).
\]
The quotient is a complete intersection with socle degree $2m+1$ and socle character zero.
\end{lemma}
\begin{proof}
For a homogeneous functional $\eta$ on $S_D$, its apolar ideal has
degree-$a$ part
$\{g\in S_a:\eta(gS_{D-a})=0\}$.
The apolar ideal of $\theta_0$ is $(A,B,z^2)$.
These quadrics annihilate $\theta_0$ and form a complete intersection
of socle degree three, on which $\theta_0$ is nonzero.

The apolar ideal of $\theta_0\circ\mathcal C_m$ is
$(A^m,B^m,z^{2m})$. Coefficient extraction gives containment, since
$A,B$ have coefficients in $\mathbb F_2$.
Equality follows from the complete-intersection socle pairing:
the functional is nonzero on $x^{4m-1}y^{m-1}z^{m-1}$.
Multiplication of the functional by $f^{m-1}$ replaces its apolar
ideal by the stated colon.

$S$ is free over $k[A,B,z]$ with basis $1,x,y,xy$.
Indeed, $x^2=A+yz$ and $y^2=B+xz$ give a monic Gr\"obner basis
with leading monomials $x^2,y^2$ for an order assigning weight two
to $x,y$ and weight one to $z$. $A,B$ belong to the coefficient ring.

Write $m=2M$. In $k[A,B,z]/(A^{2M},B^{2M},z^{4M})$,
group terms by their $z$-exponents modulo four. Put $Y=z^4$.
Since $\alpha$ is a unit, we may replace $Y$ by
$\varpi=AB+\alpha Y$. On each residue summand the algebra is
\[
\Delta[\varpi]/(\varpi^M-(AB)^M),\qquad \Delta=k[A,B]/(A^{2M},B^{2M}).
\]
It is free over $\Delta$ on $1,\varpi,\ldots,\varpi^{M-1}$. Multiplication by
$\varpi^{2M-1}=(AB)^M\varpi^{M-1}$ kills every summand with positive $\varpi$
exponent, because $\varpi^{2M}=(AB)^{2M}=0$. On the constant summand it kills
precisely the annihilator of $(AB)^M$ in $\Delta$, namely $(A^M,B^M)\Delta$. Thus its
annihilator is $(A^M,B^M,\varpi)$. Restoring the four residue summands and the free
extension gives the claimed colon equality.

The three homogeneous generators $A^M,B^M,f$ have only the origin as a common zero:
$A=B=0$ and $f=0$ force $z=0$, then $x=y=0$. Their degrees are $m,m,4$, so the socle
degree is $2m+1$. Their weights sum to $2M+M+0=0$ modulo three, as do the variable
weights, proving the socle-character assertion.
\end{proof}

\begin{lemma}\label{lem:residual-pairing}
Let $M:U\to V$ be a map of finite free $k\llbracket t\rrbracket$-modules,
and let $N\geq1$. Let $U_N\subseteq U/tU$ consist of reductions of
vectors $u$ with $Mu\in t^NV$, and define $D_N\subseteq V^*/tV^*$
in the same way for $M^*:V^*\to U^*$.
The pairing
\[
D_N\times U_N\longrightarrow k,\qquad
(v_0,u_0)\longmapsto
\langle v_0,u_0\rangle_N=[t^N]\langle\widetilde v,M\widetilde u\rangle
\]
is independent of the lifts. Its rank is the number of Smith exponents
of $M$ equal to $N$, and its left and right radicals are
$D_{N+1}$ and $U_{N+1}$.
\end{lemma}
\begin{proof}
Changing a primal lift by $ta$ changes the expression only in order
at least $N+1$, because $M^*\widetilde v$ is divisible by $t^N$.
The dual change is the same.
In Smith coordinates, a diagonal entry $t^a$ contributes to both
reduction spaces exactly when $a\geq N$, and to the pairing exactly
when $a=N$. Zero diagonal entries contribute only to the radicals.
This proves the rank and radical assertions.
\end{proof}
For multiplication by $f_t$, the socle pairing gives
\[
\langle v_0,u_0\rangle_N
=[t^N]\Lambda_q(f_t\widetilde v\widetilde u).
\]

\begin{proposition}
Fix a critical dyadic $q_0$, and suppose there are homogeneous $F,F_1$ of degree
$3q_0/2-5$ such that, modulo $(x^{q_0},y^{q_0},z^{q_0})$,
\[
\begin{gathered}
fF=0,\qquad fF_1=z^4F,\qquad \Theta=z^4F_1,\\
\Lambda_{q_0}(\Theta Fp)=\gamma\theta_0(p)\quad(p\in S_3),\quad\gamma\ne0.
\end{gathered}
\]
For $q=q_0m$, set $\widetilde u_q=f^{m-1}F^m$. By Theorem~\ref{thm:critical}, the
central primal kernel and the dual obstruction space are respectively
\[
\widetilde u_q(S/(f))_{m+1,1},\qquad
\widetilde u_q(S/(f))_{m,2},
\]
and there are $N_m=\lfloor4m/3\rfloor$ positive local factors. For dyadic $m\geq2$,
exactly $N_m-1$ positive exponents are $2m$ and precisely one is greater than $2m$.
\end{proposition}
\begin{proof}
The polynomial $\widetilde U_q=f_t^{m-1}(F+tF_1)^m$ satisfies
$f_t\widetilde U_q=t^{2m}\Theta^m$ modulo $(x^q,y^q,z^q)$, by Frobenius. These lifts
exhaust the reduction spaces in Theorem~\ref{thm:critical}, so every positive
exponent is at least $2m$. For $a$ of degree $m+1$, weight one, and $b$ of degree
$m$, weight two, the order-$2m$ pairing is
\[
\Lambda_q\bigl(f^{m-1}(\Theta F)^mab\bigr)=\gamma^m\eta_m(ab).
\]
To determine the coefficient $\gamma^m$, apply (C) with $p=\Theta F$ and
$\Psi=f^{m-1}ab$. The cubic functional $p'\mapsto\Lambda_{q_0}((\Theta F)^{\langle
m\rangle}p')$ has coefficient vector $\gamma^m\theta_0$, since $\theta_0$ has
coefficients in $\mathbb F_2$.

By the colon lemma and the complete-intersection perfect pairing, the rank on these
two character spaces is
\[
\dim_k(S/(A^{m/2},B^{m/2},f))_{m,2}=N_m-1.
\]
Indeed the two additional generators have degree $m$, opposite nonzero characters,
and are nonzero modulo $f$. Exactly one removes a dimension in character two. The
residual-pairing principle proves the assertion.
\end{proof}
\section{The exceptional elementary divisor}\label{sec:exceptional}
The first residual pairing determines all but one positive local
exponent. We compute the remaining exponent at the doubled critical
level. An integral Laurent presentation reduces the calculation to
two square blocks and one rectangular block. The square-block
determinants give the two exponents $4$ and $8$. Frobenius then
propagates the exceptional exponent to all higher levels.

\subsection{Auxiliary determinants}
For $e\geq1$, put $\ell_e=4^e$ and $\kappa'_e=(4^e-1)/3$.
Let $\Sigma_e$ be the upper shift matrix of size $\kappa'_e$, and let
$\mathsf J_e$ reverse its basis. Set
\[
\begin{gathered}\mathsf D_e=\Id+\sum_{i=1}^{e-1}\Sigma_e^{(4^e-4^i)/3},\quad
\mathsf B_e=\mathsf J_e\diag(\beta^2,\beta^5,\ldots,\beta^{4^e-2}),\\
\mathcal G_e(\beta,b)=\det(\mathsf D_e+b\mathsf B_e).\end{gathered}
\]
Here $\beta,b$ are independent indeterminates over a field of
characteristic two. We use the escape polynomials $P_j,Q_j$ from
Section~\ref{sec:escape}, evaluated at $\beta$.
\begin{theorem}
If $P_e(\alpha)=0$, then $\alpha\ne0$ and, with $\beta=\alpha+t$,
\[
\ord_t\mathcal G_e(\alpha+t,1)=2.
\]
This holds for every $e\geq1$ and every root of $P_e$ over an algebraic closure.
\end{theorem}
\begin{proof}
By Proposition~\ref{prop:escape}, the polynomials $P_j$ are squarefree,
pairwise coprime, and have no zero root. The recurrence for $Q_j$ also
shows that $Q_j$ is a product of powers of $P_i$ with $i<j$.

For $e\geq2$, write $\ell'=4^{e-1}$ and
$r=(\ell'-1)/3=\kappa'_{e-1}$, so $\kappa'_e=4r+1$.
Partition the indices into consecutive groups of sizes $r,2r+1,r$.
Every nonidentity shift in $\mathsf D_e$ is at least $\ell'=3r+1$,
so it joins only the first and last groups. Its upper-right block is
$\mathsf D_{e-1}$, since
\[
\frac{4\ell'-4^i}{3}-\ell'=\frac{\ell'-4^i}{3},
\]
and the last shift, $i=e-1$, supplies the identity in that block.

The matrix $\mathsf B_e$ preserves the middle group and exchanges the
outer groups. Put $\mathsf B'=\mathsf B_{e-1}$. Its upper-right and
lower-left blocks are $\beta^{3\ell'}\mathsf B'$ and $\mathsf B'$.
Indeed, the exponent in column $j$ is $2+3j$, and moving that column to
$3r+1+j$ adds $3\ell'$. Also
$(\mathsf B')^2=\beta^{\ell'}\Id$.

In the middle group every opposite pair of entries of $\mathsf B_e$ has product $\beta^{4\ell'}$. There are $r$ opposite pairs and one middle entry $\beta^{2\ell'}$. Its contribution to the determinant is consequently
\[
(1+b^2\beta^{4\ell'})^r(1+b\beta^{2\ell'})
=(1+b\beta^{2\ell'})^{2r+1}.
\]
After separating this middle group, the outer matrix is
\[
\begin{pmatrix}\Id&\mathsf D_{e-1}+b\beta^{3\ell'}\mathsf B'\\b\mathsf B'&\Id\end{pmatrix}.
\]
Its Schur complement is $(1+b^2\beta^{4\ell'})\Id+b\mathsf B'\mathsf D_{e-1}$. Factoring $b\mathsf B'$ and using $\det \mathsf B'=\beta^{\ell'r/2}$ gives the rational-function identity
\[
\mathcal G_e(\beta,b)=
(1+b\beta^{2\ell'})^{2r+1}b^r\beta^{\ell'r/2}
\mathcal G_{e-1}\!\left(\beta,\frac{1+b^2\beta^{4\ell'}}{b\beta^{\ell'}}\right).
\tag{*}
\]
Here $\ell'r/2$ is an integer. The identity is taken in $k(\beta,b)$.
Its left side is polynomial. The base case is
$\mathcal G_1(\beta,b)=1+b\beta^2$.

To evaluate (*) at $\beta=\alpha+t$ with $P_e(\alpha)=0$, work in the fraction field of the local ring and put
\[
\tau_0=\beta,\qquad \tau_{i+1}=\beta^3(\tau_i^4+\tau_i^{-4}).
\]
At stage $j=e-i$ of the recursion, the second argument is
\[
b_j=(\tau_i/\beta)^{2\cdot4^{j-1}}.
\]
The initial value is $b_e=1$. Substitution into (*) gives the next value:
\[
\frac{1+\tau_i^{4^j}}{\tau_i^{2\cdot4^{j-1}}\beta^{-4^{j-1}}}
=\bigl(\beta(\tau_i^2+\tau_i^{-2})\bigr)^{4^{j-1}}
=(\tau_{i+1}/\beta)^{2\cdot4^{j-2}}.
\]
For $i\geq1$, direct induction using the escape recurrences gives
\[
\tau_i=\frac{P_i^8}{\beta Q_i^8},\qquad
1+\tau_i=\frac{P_{i+1}}{\beta Q_i^8}.
\]
The initial identity $1+\tau_0=P_1$ treats $i=0$.

At the chosen root of $P_e$, every $\tau_i$, $0\leq i\leq e-1$, is a unit.
For $i\leq e-2$, the element $1+\tau_i$ is also a unit. Thus all
nonterminal factors in (*) are units, because
\[
1+b_j\beta^{2\cdot4^{j-1}}=(1+\tau_i)^{2\cdot4^{j-1}}.
\]
The only nonunit contribution is the final base determinant
\[
\mathcal G_1(\beta,b_1)=1+\tau_{e-1}^2=(1+\tau_{e-1})^2.
\]
For $e\geq2$, the expression
$1+\tau_{e-1}=P_e/(\beta Q_{e-1}^8)$ has order one.
For $e=1$, use $1+\tau_0=1+\beta$. The determinant therefore has order two.
\end{proof}

\begin{proposition}
Let $P_e(\alpha)=0$, work over $\mathcal V=k\llbracket t\rrbracket$,
and set $\beta=\alpha+t$. Put
\[
\mathsf M=\begin{pmatrix}\mathsf D_e&\mathsf B_e\\ \mathsf B_e&\mathsf D_e\end{pmatrix},\qquad
\mathsf T=\begin{pmatrix}\mathsf M&\ast\\0&\mathsf M\end{pmatrix}
\]
for any compatible block $\ast$ over $\mathcal V$. Then
$\ord_t\det \mathsf M=4$ and $\ord_t\det \mathsf T=8$.
\end{proposition}
\begin{proof}
Adding the two block rows and then the two block columns makes
$\mathsf M$ block triangular with diagonal blocks
$\mathsf D_e+\mathsf B_e$. Hence
$\det\mathsf M=\mathcal G_e(\beta,1)^2$ and
$\det\mathsf T=(\det\mathsf M)^2$. Apply the preceding theorem.
\end{proof}

\subsection{An integral Laurent presentation}
For the presentation below, $\beta$ may be any unit in
$\mathcal V=k\llbracket t\rrbracket$. Put
\[
\ell=4^e,
\quad \kappa'=(\ell-1)/3,\quad q=4\ell,\quad
\mathcal S=\mathcal V[x,y,z]/(h+\beta z^4).
\]
Weights remain $\wt(x,y,z)=(1,2,0)$. On $z=1$ set $A=x^2+y$, $B=y^2+x$, so $AB=\beta$. Write $X=1,x,y,xy$ for the four coordinate functions, with weights $0,1,2,0$.

\begin{lemma}
For $n=2\ell$ or $6\ell$, a $\mathcal V$-basis of $\mathcal S_n$ is
\[
\begin{gathered} A^r\quad(-n/2\leq r\leq n/2),\\
XA^r\quad(-n/2+1\leq r\leq n/2-1),\quad X\in\{x,y,xy\}.\end{gathered}
\]
Negative powers mean $A^{-a}=\beta^{-a}B^a$ before homogenization.
\end{lemma}
\begin{proof}
Each function in the list homogenizes to degree $n$: for $r\geq0$ multiply the homogenized $A^r=(x^2+yz)^r$ by the needed nonnegative power of $z$, and for $r<0$ use $\beta^{r}B^{-r}$ with $B=y^2+xz$. The special fibre is a smooth plane quartic because the residue of $\beta$ is nonzero. Its two points at infinity have pole orders
\[
A:(4,-4),\quad x:(2,1),\quad y:(1,2),\quad xy:(3,3).
\]
These follow from the curve equation in the charts $x=1$ and $y=1$. In the first chart, for example, $s=y/x$ is a uniformizer and $z=s^2$ times a unit. For each coordinate function, its largest Laurent exponent gives its largest pole at the first point. The four pole residues modulo four are distinct. The proposed functions are therefore linearly independent on the special fibre. Their number is $4n-2$, the rank of $\mathcal S_n$ for $n\geq4$. The latter module is free, since the quartic has unit coefficient at $z^4$ and division in $z$ gives a free normal-form basis over $\mathcal V[x,y]$. Nakayama's lemma now proves the assertion. Weight sublists give bases of the respective character summands.
\end{proof}

The cokernel of the dual central map is
\[
\mathcal S_{6\ell,2}/(x^q\mathcal S_{2\ell,1}+y^q\mathcal S_{2\ell,0}+z^q\mathcal S_{2\ell,2}).
\]
On $z=1$, multiplication by $z^q$ identifies the degree-$2\ell$ basis
with the same affine functions in degree $6\ell$. These columns are unit
pivots. Removing them leaves the Laurent coefficients outside the
degree-$2\ell$ intervals in the lemma.

There are exact identities
\[
x^q=x+W,\quad y^q=y+Z,\quad
W=\sum_{i=0}^e(A^{2\cdot4^i}+B^{4^i}),\quad Z=W(B,A).
\]
They follow by iterating $x^4+x=A^2+B$ and its conjugate. Write
$a=a_0+a_1x+a_2y+a_3xy$ and $b=b_0+b_1x+b_2y+b_3xy$.
The coordinate expressions for $x^qa+y^qb$ are
\begin{align*}
c_0&=Wa_0+Zb_0+Aa_1+Ba_3+Bb_2+Ab_3,\\
c_x&=Wa_1+Zb_1+a_0+a_3+b_2+Bb_3,\\
c_y&=Wa_2+Zb_2+a_1+Aa_3+b_0+b_3,\\
c_{xy}&=Wa_3+Zb_3+a_2+b_1.
\end{align*}
These are obtained directly from $x^2=A+y$, $y^2=B+x$.

\begin{proposition}
Let $\mathsf D=\mathsf D_e$, $\mathsf B=\mathsf B_e$ as in the preceding determinant theorem, and put
\[
\mathsf M=\begin{pmatrix}\mathsf D&\mathsf B\\ \mathsf B&\mathsf D\end{pmatrix},\quad
\mathbf E=\binom0{\mathbf e_*},\quad \mathbf G=\binom0{\mathsf J\mathbf e_*},\quad
 \mathbf e_*=\sum_{i=0}^{e-1}\mathbf e_{(4^i-1)/3}.
\]
Here standard-basis indices start at zero. The dual central cokernel is isomorphic to
\[
\coker \mathsf M\ \oplus\ \coker \mathsf T\ \oplus\ \coker \mathsf Y_e,
\qquad
\mathsf T=\begin{pmatrix}\mathsf M&\beta^{-2}\mathbf E\mathbf G^{\mathsf T}\\0&\mathsf M\end{pmatrix},
\]
where $\mathsf Y_e$ has $2\kappa'+2$ rows and $2\kappa'+1$ columns. Every elimination and rescaling used in this presentation is invertible over $\mathcal V$.
\end{proposition}
\begin{proof}
For clarity, every entry used below can be recovered by one Laurent
coefficient. Write $e_0=1,e_1=x,e_2=y,e_3=xy$, and let
$\pi_i$ extract the $e_i$ coordinate. With all columns written in
powers of $A$, the entry in row $e_iA^r$ of column $e_jA^s$ is
\[
 [A^{r-s}]\bigl(\delta_{ij}W+\pi_i(xe_j)\bigr)
 \quad\text{or}\quad
 [A^{r-s}]\bigl(\delta_{ij}Z+\pi_i(ye_j)\bigr)
\]
for an $a$ or $b$ column, respectively. A column expressed with $B^s$
has exponent $-s$ and the additional unit factor $\beta^s$.
The four coordinate identities above list all lower-coordinate terms.

First remove the columns of $a_i$ with positive Laurent exponent, using
the corresponding rows of $A^{2\ell}a_i$ above exponent $2\ell$.
Each pivot has coefficient one. The other positive powers of $W$ are at
most $\ell/2$, and the largest positive power of $Z$ is $\ell$. Neither
can reach these rows from the remaining source intervals. The
lower-coordinate terms have exponent at most $\ell+1$ and also do not
reach them. The same argument with $A,B$ exchanged removes all negative
Laurent powers of $b_i$ through the term $B^{2\ell}b_i$.
The character conditions leave the following powers:
\[
\begin{array}{c|c|c}
i&a_i\text{ in powers of }B&b_i\text{ in powers of }A\\\hline
0&1,4,\ldots,\ell&0,3,\ldots,\ell-1\\
1&0,3,\ldots,\ell-1&1,4,\ldots,\ell-3\\
2&2,5,\ldots,\ell-2&2,5,\ldots,\ell-2\\
3&1,4,\ldots,\ell-3&0,3,\ldots,\ell-1.
\end{array}
\]
The $z^q$ pivots remove scalar exponents in $[-\ell,\ell]$
and exponents in $[1-\ell,\ell-1]$ in each of $x,y,xy$.
In the scalar row, multiplication of $a_1,a_3,b_2,b_3$ by
$A,B,B,A$, respectively, shifts each exponent by one in absolute
value and leaves it in the scalar interval. In $c_y$ and $c_{xy}$
the listed powers likewise stay in the corresponding interval.
In $c_x$ the term $a_0$ reaches exponent $-\ell$.
Thus the only surviving lower-coordinate term is
$\eta B^\ell$, where $\eta$ is the coefficient of $B^\ell$ in $a_0$.
Thus only the scalar and $x$ coordinates remain coupled. At this point
there are $8\kappa'+5$ rows and $8\kappa'+4$ columns.

In the $y$ coordinate the columns are indexed by $b_{2,2+3j}$ and $a_{2,2+3j}$. Index the rows by
$A^{\ell+2+3r},B^{\ell+2+3r}$, with $0\leq j,r<\kappa'$.
The terms with matching signs give the identity and the shifts
$(\ell-4^i)/3$ in $\mathsf D$. For the other entries, use
\[
A^{2\ell}B^{2+3j}
 =\beta^{2+3j}A^{2\ell-2-3j}.
\]
This reaches row $r$ exactly when $j=\kappa'-1-r$.
The conjugate identity gives the other off-diagonal block. The $y$
coordinate therefore gives $\mathsf M$.

In the scalar coordinate, the row at $B^{2\ell}$ imposes $b_{0,0}=\eta$, a unit pivot. For the remaining scalar columns put
\[
p_j=b_{0,3+3j},\qquad r_j=\beta^{-1}a_{0,1+3j},\qquad v_*=(p,r).
\]
Here and below $0\leq j<\kappa'$. Use the rows
$A^{\ell+3+3r}$ and $B^{\ell+1+3r}$, $0\leq r<\kappa'$, dividing the
second group by $\beta$. For these choices the within-sign entries, in both row groups, are
\[
D_{rj}=[r=j]+\sum_{i=1}^{e-1}[j-r=(\ell-4^i)/3],
\]
and the cross entries are
\[
B_{rj}=\beta^{2+3j}[r+j=\kappa'-1].
\]
The residual boundary column has entry $\beta^{-1}$ in the second
row group at $r=(4^i-1)/3$, for $0\leq i<e$, and is zero elsewhere.
Consequently the scalar equations are
\[
\mathsf Mv_*+\beta^{-1}\mathbf E\eta=0.
\]
Indeed, the cross powers are $\beta^{1+3j}$ before the column scaling, so they become $\beta^{2+3j}$ in the positive rows. The negative rows give the same entries after row scaling. The only additional high terms are $\eta B^{\ell+4^i}$, $0\leq i<e$, whose row indices are $(4^i-1)/3$.

In the $x$ coordinate, the row at $A^{2\ell}$ eliminates $a_{1,0}$ by a unit pivot. Put
\[
s_j=b_{1,1+3j},\qquad t_j=\beta a_{1,3+3j},\qquad w_*=(s,t).
\]
Use rows $A^{\ell+1+3r}$ and $B^{\ell+3+3r}$,
$0\leq r<\kappa'$, multiplying the second group by $\beta$.
Their equations are $\mathsf Mw_*=0$. The remaining boundary row,
at $B^\ell$, is
\[
\mathbf G^{\mathsf T}w_*+\beta\eta=0,
\]
since its contributing $a_1$ indices satisfy
\[
\frac{\ell-4^i}{3}-1=\kappa'-1-\frac{4^i-1}{3},\qquad 0\leq i<e.
\]
Eliminating $\eta$ with this unit coefficient $\beta$ gives the asserted $4\kappa'\times4\kappa'$ block $\mathsf T$ on $(v_*,w_*)$.

In the $xy$ coordinate only the coefficients of $Wa_3+Zb_3$
remain. The $a_3,b_3$ columns have no entries in the other surviving
rows. They therefore give a direct summand $\mathsf Y_e$ with
$2\kappa'+2$ rows and $2\kappa'+1$ columns. Indeed, the surviving columns correspond to $A^{3j}$ in $b_3$
for $0\leq j\leq\kappa'$ and to $B^{1+3j}$ in $a_3$ for
$0\leq j<\kappa'$. The rows correspond to $A^{\ell+3r}$ and
$B^{\ell+1+3r}$ for $0\leq r\leq\kappa'$.

The unit eliminations involving the scalar and $x$ coordinates can
be recorded as follows:
\[
\begin{array}{c|c|c|c}
\text{row}&\text{eliminated variable}&\text{pivot}&\text{substitution}\\\hline
B^{2\ell}\text{ in }c_0&b_{0,0}&1&b_{0,0}=\eta\\
A^{2\ell}\text{ in }c_x&a_{1,0}&1&\text{solve its row equation}\\
B^{\ell}\text{ in }c_x&\eta&\beta&
\eta=\beta^{-1}\mathbf G^{\mathsf T}w_*
\end{array}
\]
The last substitution in
$\mathsf Mv_*+\beta^{-1}\mathbf E\eta=0$ gives
$\mathsf Mv_*+\beta^{-2}\mathbf E\mathbf G^{\mathsf T}w_*=0$.
Together with $\mathsf Mw_*=0$, this gives $\mathsf T$.

The three boundary pivots remove three rows and columns. The remaining
blocks thus have $8\kappa'+2$ rows and $8\kappa'+1$ columns in total,
as does the direct sum in the statement. All pivots are units, and all
rescalings use powers of $\beta$.
\end{proof}

\begin{proposition}\label{prop:doubled}
Let $P_e(\alpha)=0$ and set $\beta=\alpha+t$ in the preceding
presentation. At $q=4^{e+1}$, the dual central map has free cokernel
rank one and exactly two positive local exponents. These exponents are
$4$ and $8$.
\end{proposition}
\begin{proof}
Proposition~\ref{prop:rank} gives free cokernel rank one.
Theorem~\ref{thm:critical}, with $m=2$, gives exactly two positive
exponents. The two square blocks have determinant orders $4$ and $8$.
Each therefore contributes at least one positive exponent. It follows
that each contributes exactly one and that $\mathsf Y_e$ contributes none.
The two exponents equal the corresponding determinant orders.
\end{proof}

\subsection{The order-eight pairing}
Fix a root $\alpha$ of $P_e$, put $q_0=2\cdot4^e$, and retain the
critical $F_0,F_1$ from Section~\ref{sec:obstruction}. Write $F=F_0$,
$\xi=fF^2$ and $q_1=2q_0$.
\begin{lemma}\label{lem:order-eight}
There are lifts $V_A,V_B$ of $\xi A,\xi B$ modulo $t^8$ such that
$f_tV_D=t^8\Theta_D$ modulo $(x^{q_1},y^{q_1},z^{q_1})$, and
\[
\Lambda_{q_1}(\Theta_A\xi Ay)\ne0,\qquad
\Lambda_{q_1}(\Theta_B\xi Bx)\ne0.
\]
Write $\lambda_A=\Lambda_{q_1}(\Theta_A\xi Ay)$ and $\lambda_B=\Lambda_{q_1}(\Theta_B\xi Bx)$.
The $V_D$ may be chosen of degree at most seven in $t$.
\end{lemma}
\begin{proof}
The apolar calculation at $m=2$ gives the dual radical
$\xi\langle A,B\rangle$ and the primal radical $\xi(AS_1+BS_1)$ of
$\langle\,,\rangle_4$. Proposition~\ref{prop:doubled} and its $x,y$ conjugate give
exponents $4,8$ in the two nonzero characters. In character zero the dual
radical is zero, so both exponents are four. In either nonzero character the dual radical at order four is one-dimensional. Its positive Smith exponent is eight, so it survives to order eight. Thus $\langle\,,\rangle_8$ has no left radical on
$\xi\langle A,B\rangle$.
Since $\xi$ has weight zero and $B,y$ have weights one and two,
$\xi By$ lies in the weight-zero part of the primal radical at order four.
Every positive exponent in this character is four. Smith normal form
therefore gives a vector $v$ in the kernel over $k\llbracket t\rrbracket$
whose reduction is $\xi By$. Translation by $g_1$ preserves $f_t$,
fixes $\xi$ and $B$, and sends $y$ to $y+z$. Hence $g_1v-v$ lies
in the kernel and reduces to $\xi Bz$. Its projection to the character
of $\xi Bz$ is a kernel lift in that character. In particular,
$\langle\xi A,\xi Bz\rangle_8=0$. The primal radical at order four
in the character opposite to $\xi A$ is spanned by $\xi Ay$ and $\xi Bz$.
Thus $\langle\xi A,\xi Ay\rangle_8\ne0$. Truncate a lift of $\xi A$ after $t^7$
to obtain $V_A$. Since $f_t$ is linear in $t$, its image is $t^8\Theta_A$
with $\Theta_A$ independent of $t$. Symmetry gives $V_B,\Theta_B$.
\end{proof}
\subsection{Frobenius propagation}
\begin{theorem}\label{thm:finite-orders}
Let $\alpha$ be a root of $P_e$ and $q_0=2\cdot4^e$.
At $q=q_0$ the sole positive central exponent is two.
At every dyadic $q=q_0m$, $m\geq2$, the positive exponents are
\[
\underbrace{2m,\ldots,2m}_{\lfloor4m/3\rfloor-1\ \text{times}},\quad4m.
\]
There are no positive exponents at this parameter for $q<q_0$.
\end{theorem}
\begin{proof}
At $q_0$ apply Theorem~\ref{thm:critical-order}. Below $q_0$ apply Theorem~\ref{thm:support}. For $m\geq2$, write $m=2r$, so $q=rq_1$, and use the lifts in Lemma~\ref{lem:order-eight}. Choose $D=A$ for $r\equiv1\pmod3$ and $D=B$ for $r\equiv2\pmod3$. Put
\[
\mathcal W_r(t)=f_t^{r-1}V_D(t)^r.
\]
Frobenius gives $f_t\mathcal W_r=t^{8r}\Theta_D^r$ modulo $(x^q,y^q,z^q)$, and
$\mathcal W_r(0)=f^{2r-1}F^{2r}D^r=\widetilde u_qD^r$.
Its degree is $3q/2-4$ and weight two. Thus $\mathcal W_r$ and $y\mathcal W_r$ give a dual and primal lift modulo $t^{8r}$ in the required central character spaces. Their residual pairing equals
\[
\Lambda_q\bigl(yf^{r-1}(\Theta_D\xi D)^r\bigr).
\]
Formula (C) and the extraction identities turn this into
\[
\begin{cases}
\Lambda_{q_1}\bigl(y(\Theta_A\xi A)^{\langle r\rangle}\bigr)=\lambda_A^r,&r\equiv1\pmod3,\\
\Lambda_{q_1}\bigl(x(\Theta_B\xi B)^{\langle r\rangle}\bigr)=\lambda_B^r,&r\equiv2\pmod3.
\end{cases}
\]
Both scalars are nonzero. The residual pairing therefore detects an
exponent $8r=4m$. The first residual pairing shows that only one exponent
exceeds $2m$, which proves the formula.
\end{proof}

For example, the finite exponents can be read off from the pencil
over $K[t]/(t^b)$, where $K$ is the residue field. If $d_b$ denotes
the $K$-dimension of its cokernel and $d_0=0$, then $d_b-d_{b-1}$
counts the exponents at least $b$. At $\alpha=1$, elimination gives
$4,8$ for $q=16$ and $8,8,8,8,16$ for $q=32$. At a root of $P_2$,
it gives a single exponent $2$ for $q=32$ and the pair $4,8$ for $q=64$.

\section{The local calculation at infinity}\label{sec:infinity}
The remaining point of the parameter line is $[u:w]=[1:0]$.
On the chart $u\ne0$, put $s=w/u$. The quartic equation becomes
$z^4+sh$, whose special fibre is the quadruple line $z^4=0$.
The smooth-curve calculation no longer applies. We instead use monic
division in $z$ to find the local length of $\mathcal T_q$ at this point.

Work over $\mathcal V=k\llbracket s\rrbracket$ in characteristic two, put $h=x^2y^2+(x^3+y^3)z+xyz^2$, and let
\[
q=4N\geq16,\qquad f=z^4+sh,\qquad
\mathcal R=\mathcal V[x,y,z]/(x^q,y^q,f),\quad j=6N+1.
\]
The integer $N$ is dyadic. Let $\NF$ denote monic division by $f$ in
$\mathcal V[x,y,z]$, before imposing $x^q=y^q=0$. Its values have $z$
exponent below four. Each reduction replaces $z^4$ by $sh$, so it never
decreases the $x$ or $y$ exponent. Normal forms in $\mathcal R$ are
therefore obtained by discarding the terms with $x$ or $y$ exponent at
least $q$.

We also use monic division over
$\mathcal V[x^{\pm1},y^{\pm1}]$. Every Laurent expression below is
computed in this ambient coefficient ring and then projected to the
specified exponent ranges.

\begin{lemma}
Let $r=\lfloor8N/3\rfloor$. The weight-one map
\[
\mathcal K_N=h^N:\mathcal R_{2N+1,1}\longrightarrow\mathcal R_{6N+1,1}
\]
has $r$ rows and $r+1$ columns, and the central cokernel at infinity is $\coker(s^N\mathcal K_N)$. Moreover, in the quotient by $(x^{2N},y^{2N},f)$,
\[
h^N=s^{N-1}(xy)^{2N-1}z^2.
\tag{B}
\]
\end{lemma}
\begin{proof}
The free monomial bases are $x^ay^bz^c$ with $a,b<4N$ and $c<4$. Counting those of the two specified degrees and character gives $r+1,r$ respectively. The relation $z^4=sh$ gives $z^q=s^Nh^N$. Factoring out the $z^q$ relation therefore yields precisely the cokernel stated.

For a dyadic $m\leq N$,
\[
h^m=x^{2m}y^{2m}+(x^{3m}+y^{3m})z^m+x^my^mz^{2m}.
\]
Multiplication by $(xy)^{2N-2m}$ kills the first three terms modulo $x^{2N},y^{2N}$. If $m\geq2$, substitute $z^{2m}=s^{m/2}h^{m/2}$ in the last term. Repeated halving gives the total factor $s^{N-1}$ and leaves $(xy)^{2N-2}h$. Its only surviving term is $(xy)^{2N-1}z^2$. This also proves the base case $N=1$ of (B).
\end{proof}

Call a source monomial good when its $x,y$ exponents are both below $2N$. Modulo $s$, multiplication by $h^N$ is multiplication by $\mathsf X=x^{2N}y^{2N}$. It identifies the good source columns with all but one target row. There are precisely two bad source columns. In the original character-one coordinates they are
\[
\begin{array}{c|ccc}
&b_0&b_1&m_*\\\hline
N\equiv1\pmod3&x^{2N}y&y^{2N}z&x^{2N-1}y^{4N-1}z^3\\
N\equiv2\pmod3&y^{2N+1}&x^{2N}z&x^{4N-1}y^{2N-1}z^3.
\end{array}
\]
Ordering the good rows to match the good columns gives
\[
\mathcal K_N=\begin{pmatrix}\mathsf H_N&u_N&v_N\\w_N&0&s^{N-1}\end{pmatrix},\qquad \mathsf H_N\equiv \Id\pmod s.
\]
The zero follows because $b_0$ already exceeds the missing row in one coordinate. (B) gives the last entry. Thus $\mathcal K_N$ is equivalent over $\mathcal V$ to an identity block and the Schur row
\[
(\rho_N,\rho'_N)=
(w_N\mathsf H_N^{-1}u_N,\ s^{N-1}+w_N\mathsf H_N^{-1}v_N).
\]

\begin{lemma}\label{lem:infty-coefficient}
For dyadic $n\geq1$ set
\[
\phi_n=\begin{cases}
[x^{4n-1}y^{8n-2}z^3]\NF(h^{3n}),&n\equiv1\pmod3,\\
[x^{4n-2}y^{8n-1}z^3]\NF(h^{3n}),&n\equiv2\pmod3.
\end{cases}
\]
Then $\phi_n=s^{n-1}$.
\end{lemma}
\begin{proof}
The coefficient of $z^3$ in a remainder modulo the monic quartic is the coefficient of $z^{-1}$ in the Laurent expansion of the original polynomial divided by $f$ at $z=\infty$. For $n\geq2$ use
\[
\frac{h^{3n}}f=f\left(\frac{h^{3n/2}}f\right)^2.
\]
In the first coefficient above, parity selects only $sx^3z$ from $f$. In the second it selects only $sy^3z$. Halving the remaining exponents exchanges the two cases. Consequently $\phi_n=s\phi_{n/2}^2$. Direct division for $n=1$ gives $\phi_1=1$, so $\phi_n=s^{n-1}$.
\end{proof}

\begin{theorem}\label{thm:infinity}
For every dyadic $N\geq4$,
\[
\rho_N\equiv s^{N-1}\pmod{s^N},\qquad
\rho'_N\equiv0\pmod{s^N}.
\]
Consequently the positive local exponents of the original central pencil at infinity are
\[
\underbrace{N,\ldots,N}_{r-1\ \text{times}},\quad 2N-1,
\]
and the order of its maximal-minor gcd at infinity is
\[
\nu_\infty(G_q)=N(r+1)-1
=\frac q4\left(\left\lfloor\frac{2q}{3}\right\rfloor+1\right)-1.
\]
\end{theorem}
\begin{proof}
The good-row block is invertible. It suffices to compute its two
Schur entries modulo $s^N$, since an entry of order $N-1$ and a
second entry divisible by $s^N$ give the last invariant factor.
Write $N=4n$ and set
\[
a=\NF(h^n),\quad b=\NF(h^{2n})=\NF(a^2),\quad
\mathcal A=(x^{3N}+y^{3N})a,\quad\mathcal B=x^Ny^Nb.
\]
All products in this proof are reduced to normal form. Frobenius gives
$h^N=\mathsf X+s^n\mathcal A+s^{2n}\mathcal B$.

Let $\Pi$ retain the normal monomials whose $x,y$ exponents both lie
in $[0,2N-1]$. On the space of good source columns, define
\[
\Delta_D(g)=\Pi\bigl(\mathsf X^{-1}\NF(Dg)\bigr),\qquad
\lambda_D(g)=[m_*]\NF(Dg),\qquad D\in\{\mathcal A,\mathcal B\}.
\]
Multiplication by $\mathsf X^{-1}$ here means shifting each $x,y$
exponent down by $2N$. A term discarded by the truncation at $q=4N$
still has an exponent at least $2N$ after this shift and is removed by
$\Pi$. Thus these operators give the entries of the truncated matrix.
We obtain
\[
\Delta:=\mathsf H_N-\Id
 =s^n\Delta_{\mathcal A}+s^{2n}\Delta_{\mathcal B},\qquad
w_N=s^n\lambda_{\mathcal A}+s^{2n}\lambda_{\mathcal B}.
\]
The same expressions applied to $b_0,b_1$ give $u_N,v_N$.
All entries of $\Delta,w_N,u_N,v_N$ are divisible by $s^n$.
Hence the expansion
$\mathsf H_N^{-1}=\Id+\Delta+\Delta^2+\cdots$ gives, modulo $s^{4n}$,
only $w_Nu_N+w_N\Delta u_N$ and the analogous expression with $v_N$.
Terms with four perturbation factors have order at least $4n$.
Normal division introduces only nonnegative powers of $s$.
The table records the terms below order $s^{4n}$, with the common
power of $s$ in the second column removed. Operators act from right to
left. The projections $a'$ and $\Pi'$ and the coefficient functional
$\Lambda'_\epsilon$ are specified below.
\[
\begin{array}{c|c|c|c}
\text{product}&\text{power}& b_0 & b_1\\ \hline
\mathcal A^2&2n&0&s^{2n-1}\\
\mathcal A\mathcal B&3n&\Lambda'_\epsilon(a\Pi'(b))&0\\
\mathcal B\mathcal A&3n&\Lambda'_\epsilon(ba')&0\\
\mathcal A^3&3n&\Lambda'_\epsilon(a\Pi'(aa'))&0
\end{array}
\]
The $b_1$ column also has the direct boundary term $s^{N-1}$.
Two $\mathcal B$ factors already have order $4n$.

Exchange $x,y$ if necessary so that the missing row is always
$m_*=x^{2N-1}y^{4N-1}z^3$. In these coordinates $b_0=x^{2N}\epsilon$ with $\epsilon=y$ if $N\equiv1\pmod3$ and $\epsilon=x$ if $N\equiv2\pmod3$. $b_1=y^{2N}z$. The exchange changes the selected character but preserves all normal-form and exponent arguments.

\emph{The first Schur entry.}
Define $a'$ by $\Pi(\epsilon y^Na)=\epsilon y^Na'$. Every $x,y$ exponent in $a$ is at most $N-1$: its degree is $N$, and its $z$-constant terms are divisible by $x^2y^2$. Setting $x=0$ gives the exact expression
\[
a(0,y,z)=\begin{cases}
s^{(n-1)/3}y^{N-1}z,&n\equiv1\pmod3,\\
s^{(n-2)/3}y^{N-2}z^2,&n\equiv2\pmod3.
\end{cases}
\]
Indeed $h^n=y^{3n}z^n$ there, and each use of $z^4=sy^3z$ lowers the $z$ exponent by three. Since $N\equiv n\pmod3$, we have $a'=a+\partial$, where
$\partial=s^{(n-1)/3}y^{N-1}z$ when $\epsilon=y$, and $\partial=0$ when $\epsilon=x$.

Let $\Pi'$ be the projection characterized by
$\Pi(\epsilon x^Ny^{-N}g)=\epsilon x^Ny^{-N}\Pi'(g)$. Its allowed ranges are
\[
\begin{array}{c|cc}
&x\text{ exponent}&y\text{ exponent}\\\hline
\epsilon=y&[0,N-1]&[N-1,3N-2]\\
\epsilon=x&[0,N-2]&[N,3N-1].
\end{array}
\]
Let $\Lambda'_\epsilon(g)=[m_*](\epsilon x^Ny^{2N}g)$. The two initial good-column vectors are
\[
u_{\mathcal A}=\epsilon y^Na',\qquad
u_{\mathcal B}=\epsilon x^Ny^{-N}\Pi'(b).
\]
The $x^{3N}$ branch cannot enter the missing row at the last step. In $\lambda_{\mathcal A}u_{\mathcal A}$, the remaining branch has $y$ exponent at least $4N$, so this two-perturbation term vanishes. The two mixed terms contribute
$s^{3n}\Lambda'_\epsilon(a\Pi'(b)+ba')$.
The only three-perturbation term below $s^{4n}$ is
$s^{3n}\lambda_{\mathcal A}\Delta_{\mathcal A}u_{\mathcal A}$. Its middle $y^{3N}$ branch is killed by $\Pi$, while its other branch gives
$\Delta_{\mathcal A}u_{\mathcal A}=\epsilon x^Ny^{-N}\Pi'(aa')$.
Thus
\[
\rho_N\equiv s^{3n}\Lambda'_\epsilon\bigl(ba'+a\Pi'(a\partial)\bigr)\pmod{s^{4n}},
\]
using $aa'=b+a\partial$ to cancel the two $a\Pi'(b)$ terms.

If $\partial\ne0$, then $\Pi'(a\partial)=a\partial$. Its $y$ exponents are at least $N-1$. Its $x$ exponents are below $N$: in $\NF(za)$, an $x$ exponent at least $N$ would leave total $y,z$ degree at most one. The $z$-constant part is divisible by $x^2y^2$, and at $y=0$ the normal form is $s^{(n-1)/3}x^{N-1}z^2$, excluding all remaining possibilities. The upper $y$ bound follows from $\deg(a\partial)=2N\leq3N-2$. Therefore in both cases the expression inside $\Lambda'_\epsilon$ is $ab$, and
\[
\rho_N\equiv s^{3n}\phi_n\pmod{s^{4n}},
\]
with $\phi_n=s^{n-1}$ by Lemma~\ref{lem:infty-coefficient}. This proves $\rho_N\equiv s^{N-1}\pmod{s^N}$.

The first entry already gives the bound needed for Theorem A. Indeed,
\[
 \nu_\infty(G_q)=Nr+\min\{\ord_s\rho_N,\ord_s\rho'_N\}
                 \leq Nr+N-1 .
\]
Since $c_q=n_q-\deg G_q$, this upper bound and the finite orders
give Proposition~\ref{prop:index-lower}. The second entry gives equality.

\emph{The second Schur entry.}
Put $c=\NF(za)$ and let $c'$ retain its terms with $x$ exponent below $N$. Let $\Pi''$ be the projection retaining normal monomials whose $x$ exponent is at least $N$ and $y$ exponent below $N$. The two initial vectors for $b_1$ are
\[
v_{\mathcal A}=x^Nc',\qquad
v_{\mathcal B}=x^{-N}y^N\Pi''(\NF(zb)).
\]
No further upper $x$ bound is needed by the total degree. The two-$\mathcal A$ contribution is
\[
\lambda_{\mathcal A}v_{\mathcal A}
=[x^{N-1}y^{N-1}z^3](ac')
=[x^{N-1}y^{N-1}z^3]\NF(zh^{2n})
=s^{2n-1},
\]
by (B) with $2n$ in place of $N$. Deleted terms of $c$ cannot reach the specified $x$ exponent, since normal division never decreases it. The mixed terms vanish: the final $\mathcal A$ acting on $v_{\mathcal B}$ has $y$ exponent at least $4N$, and the final $\mathcal B$ acting on $v_{\mathcal A}$ has $x$ exponent at least $2N$.

The only possible three-perturbation term is $\lambda_{\mathcal A}\Delta_{\mathcal A}v_{\mathcal A}$. In its middle step the $x^{3N}$ branch is killed by $\Pi$ because its $x$ exponent after division by $\mathsf X$ is already at least $2N$. The other branch is $x^{-N}y^N\Pi''(ac')$. Applying the final $y^{3N}$ branch overshoots the missing row's $y$ exponent, and the final $x^{3N}$ branch always overshoots its $x$ exponent. Thus that term vanishes as well. Including the direct boundary entry proves
\[
\rho'_N\equiv s^{N-1}+s^{2n}s^{2n-1}=0\pmod{s^N}.
\]

Finally $\mathsf H_N$ is invertible over $\mathcal V$, and the Schur row has gcd valuation exactly $N-1$ by the first congruence. Thus $\mathcal K_N$ has $r-1$ unit factors and one factor of valuation $N-1$. Multiplication by $s^N$ adds $N$ to all $r$ exponents, proving both the list and the gcd order.
\end{proof}

The two entries can also be computed in the original monomial bases
over $\mathbb F_2[s]/(s^N)$. Reduce each column of multiplication by
$h^N$ using $z^4=sh$ and $x^{4N}=y^{4N}=0$, then eliminate the good
columns. For $N=4,8,16,32,64$, the resulting row is $(s^{N-1},0)$.

\section{Determinantal factorization and the main results}\label{sec:main}

The preceding sections determine the support and every local length
of the torsion sheaf $\mathcal T_q$ in~\eqref{eq:kernel-sheaf}.
Their sum is $\deg G_q$. Formula~\eqref{eq:index-gcd} then gives $c_q$,
and local duality yields the regularity bounds.

\subsection{The minimal index}

\begin{proposition}\label{prop:index-lower}
For every dyadic $q\geq4$, the coefficient degree $c_q$ of the
primitive kernel vector satisfies
\[
 c_q\geq b_q.
\]
For $q\geq16$ this bound uses the finite local orders and only the
first Schur entry at infinity.
\end{proposition}
\begin{proof}
Proposition~\ref{prop:rank} proves generic full row rank and a one-dimensional generic kernel. Hence the cofactor lemma gives
$c_q=q^2/4-2-\deg G_q$.

\emph{Local orders.}
Theorems~\ref{thm:support} and~\ref{thm:finite-orders} determine all finite orders. At a root of $P_j$ put $q_0=2\cdot4^j$. There are no factors below $q_0$. At $q=q_0$ the order is two. For $q=q_0m$, $m\geq2$, the total order is
\[
2m\left(\left\lfloor\frac{4m}{3}\right\rfloor-1\right)+4m
=2m\left(\left\lfloor\frac{4m}{3}\right\rfloor+1\right).
\]
Define on dyadic arguments
\[
a(s)=\begin{cases}0&s\leq4,\\2&s=8,\\
\dfrac s4\left(\left\lfloor\dfrac s6\right\rfloor+1\right)&s\geq16.
\end{cases}
\]
Since $\deg P_j=8^{j-1}$, the finite contribution to $\deg G_q$ is
\[
\nu_{\mathrm{fin}}(q)=\sum_{j\geq1}8^{j-1}a(q/4^{j-1}).
\]
Put
\[
 U(q)=\frac q4\left(\left\lfloor\frac{2q}{3}\right\rfloor+1\right)-1.
\]
For $q\geq16$ the first Schur entry in the proof of
Theorem~\ref{thm:infinity} gives $\nu_\infty(q)\leq U(q)$.
Thus $c_q\geq q^2/4-2-\nu_{\mathrm{fin}}(q)-U(q)$.

\emph{The initial values.}
At $q=4$, $z^4=0$ in $A_4$, so the gcd is $w^2$ and $c_4=0$. At $q=8$ the character-one vectors
\begin{align*}
v_0={}&y^2z^7+y^5z^4+xy^3z^5+x^2y^4z^3+x^3y^2z^4\\
&+x^4y^3z^2+x^5yz^3+x^6y^2z+x^7z^2,\\
v_1={}&y^2z^7+x^4z^5
\end{align*}
satisfy $hv_0=0$, $hv_1=z^4v_0\ne0$, $z^4v_1=0$. These follow by multiplying the displayed four-term $h$ and discarding monomials with an exponent at least eight. Since $\ker h$ in this character is one-dimensional, there is a degree-one pencil kernel vector but no degree-zero one. Hence $c_8=1$. Thus $\deg G_8=13$. Its only finite factor has order two, so its infinity order is eleven. Thus $\nu_\infty(q)=U(q)$ for $q=4,8$.

\emph{The recurrence.}
Put $v(q)=q^2/4-1-U(q)-\nu_{\mathrm{fin}}(q)$. Then $c_q+1\geq v(q)$.
For every dyadic $q\geq4$,
\[
\begin{aligned}
\nu_{\mathrm{fin}}(4q)&=a(4q)+8\nu_{\mathrm{fin}}(q),\\
a(4q)&=4U(q)+4,\\
U(4q)&=4U(q)+2q^2+3.
\end{aligned}
\]
The last equality uses $\lfloor8q/3\rfloor=2q+\lfloor2q/3\rfloor$.
Substitution gives $v(4q)=8v(q)$. The initial values $v(4)=1$
and $v(8)=2$ give $v(q)=b_q+1$ on the two sequences
$q=4^r$ and $q=2\cdot4^r$. Hence $c_q\geq b_q$.
\end{proof}

\begin{theorem}\label{thm:index}
For every $q=2^e\geq4$, the kernel index is
\[
 c_q=b_q=2^{e+\lfloor e/2\rfloor-3}-1.
\]
\end{theorem}
\begin{proof}
The second Schur entry in Theorem~\ref{thm:infinity} has order
at least $q/4$, so $\nu_\infty(q)=U(q)$ for $q\geq16$. Equality at
$q=4,8$ was checked above. The cofactor formula now gives
$c_q+1=v(q)=b_q+1$.
\end{proof}

A direct computation of $c_q$ uses only $H=\mu_h$ and $Z=\mu_{z^4}$.
Choose $v_0$ spanning $\ker H$ in the central character and a right
inverse $L$ of $H$. Put $v_{i+1}=LZv_i$. Then $c_q$ is the least $d$
such that $Zv_d\in\langle Zv_0,\ldots,Zv_{d-1}\rangle$, with the
empty span taken to be zero. Indeed, a relation for $H+tZ$ with
constant term $v_0$ satisfies $Ha_i=Za_{i-1}$ at each step and
$Za_d=0$ at the last. Each $a_i$ is determined up to a multiple
of $v_0$, so the last equation is exactly the stated dependence.
In the monomial bases this gives $c_q=0,1,7,15,63,127,511$ for
$q=4,8,16,32,64,128,256$, respectively.

\subsection{Regularity}
\begin{theorem}\label{thm:regularity}
Let $R$ and $I$ be as in~\eqref{eq:fixed-pair}. The ring $R$ is a
four-dimensional geometrically normal hypersurface domain, and $I$ is a
height-two prime. For every $e\geq3$,
\[
\reg_R I^{[2^e]}\geq\frac32\,2^e+2^{e+\lfloor e/2\rfloor-3}.
\]
In particular, $\reg_R I^{[2^e]}/2^e\to\infty$.
For $q=4^r$, $r\geq2$, the lower bound is $3q/2+q^{3/2}/8$.
\end{theorem}
\begin{proof}
Proposition~\ref{prop:normal} gives the assertions about $R$ and $I$.
For $q=2^e$, $e\geq3$, apply Proposition~\ref{prop:index-lower} in
Proposition~\ref{prop:regularity-reduction}.
\end{proof}
\begin{corollary}\label{cor:polynomial-regularity}
Let $P=k[x,y,z,u,w]$, $J=(uz^4+wh)$ and $K=(x,y,z)$.
Then $\reg_P(J+K^{[q]})/q\to\infty$ along powers of two.
\end{corollary}
\begin{proof}
Put $\mathfrak L_q=J+K^{[q]}$. The quotient $P/\mathfrak L_q=R/I^{[q]}$ has a nonzero
$H^0$ class in degree $3q/2+c_q$ by the proof of
Proposition~\ref{prop:regularity-reduction}. Since
$H^1_{P_+}(\mathfrak L_q)\simeq H^0_{P_+}(P/\mathfrak L_q)$,
$\reg_P \mathfrak L_q\geq3q/2+c_q+1$. Apply Proposition~\ref{prop:index-lower}.
\end{proof}

\section{Local cohomology, associated primes, and primary decompositions}\label{sec:primary}

Keep $T_q=R/I^{[q]}$, $S'=k[u,w]$, and
$\mm=(x,y,z,u,w)R$. Throughout this section $q=2^e\geq8$, so $c_q\geq b_q>0$.
Put
\[
L_q=H^0_\mm(T_q),\qquad
\lambda_q=\min\{a\geq0:\mm^aL_q=0\},\qquad
J_q=I^{[q]}:\mm^\infty.
\]
Then $L_q=J_q/I^{[q]}$.

\subsection{Local cohomology and saturation}
\begin{theorem}\label{thm:LC}
For every dyadic $q\geq8$,
\[
w^{c_q-1}L_q\ne0,\qquad
\lambda_q\geq c_q,\qquad
\dim_kL_q\geq c_q(c_q+1).
\]
The saturation exponent is
\[
\lambda_q=\min\{a\geq0:I^{[q]}:\mm^a=J_q\}.
\]
There is no fixed integer $B$ with $\mm^{Bq}L_q=0$ for every $q$.
Hence $(R,I)$ does not satisfy $(\mathrm{LC})$.
\end{theorem}
\begin{proof}
Let $\mathcal N_q$ be the first-degree $3q/2+1$, character-one component
of $T_q$, viewed as an $S'$-module. It is the central cokernel.
Lemma~\ref{lem:duality} gives
\[
H^0_{(u,w)}(\mathcal N_q)^\vee
\simeq \bigl(S'/(u,w)^{c_q}\bigr)(c_q-1).
\]
Here $(-)^\vee$ denotes graded $k$-duality. Multiplication by
$w^{c_q-1}$ is nonzero on $S'/(u,w)^{c_q}$, so it is nonzero on its
dual. Hence
\[
w^{c_q-1}H^0_{(u,w)}(\mathcal N_q)\ne0,
\qquad
\dim_kH^0_{(u,w)}(\mathcal N_q)=\binom{c_q+1}{2}.
\]
The ideal $(x,y,z)$ acts nilpotently on $T_q$. It follows that
$H^0_{(u,w)}(T_q)=H^0_\mm(T_q)$, and the displayed module is a direct
summand of $L_q$ as an $S'$-module. Interchanging $x$ and $y$ gives an
isomorphic summand in character two. The two character summands are
disjoint. This proves the nonvanishing and the length bound, and the
nonvanishing gives $\lambda_q\geq c_q$.

Since $L_q=J_q/I^{[q]}$, its annihilation by $\mm^a$ is equivalent to
$J_q\subseteq I^{[q]}:\mm^a$, hence to equality. This proves the
formula for $\lambda_q$. Finally,
$c_q/q\to\infty$ by Proposition~\ref{prop:index-lower}. Thus $\lambda_q$ cannot
be bounded by $aq$ for a fixed integer $a$.
\end{proof}

\begin{corollary}\label{cor:dimension}
Let $k$ be a field of characteristic two. For every $d\geq4$ there
exist a geometrically normal standard graded hypersurface $R_d$ of
dimension $d$ over $k$ and a homogeneous prime $I_d$ such that
$R_d/I_d\simeq k[u,w]$ and $(R_d,I_d)$ fails $(\mathrm{LC})$.
\end{corollary}
\begin{proof}
Write $s=d-4$, set $R_d=R[t_1,\ldots,t_s]$ with $\deg t_i=1$,
and put $I_d=IR_d+(t_1,\ldots,t_s)$. Polynomial extension preserves
geometric normality, and $R_d/I_d=R/I$. For $q=2^e$ let
$B_q=k[t_1,\ldots,t_s]/(t_1^q,\ldots,t_s^q)$. Then
\[
 R_d/I_d^{[q]}\simeq T_q\otimes_k B_q,\qquad
 H^0_{(R_d)_+}(R_d/I_d^{[q]})\simeq L_q\otimes_k B_q.
\]
Indeed, $(t_1,\ldots,t_s)$ acts nilpotently on the tensor product,
so torsion with respect to $(R_d)_+$ equals torsion with respect to
$\mm$. The monomial basis of $B_q$ identifies the latter with a
direct sum of copies of $L_q$. Hence $w^{b_q-1}$ acts nontrivially
for every $q\geq8$, as in Theorem A.
\end{proof}

\begin{corollary}\label{cor:vertex-saturation}
For $q\geq8$, the vertex defined by $\mm$ is an embedded point of $Z_e$, and
\[
\depth(T_q)_{\mathfrak m}=0.
\]
If
\[
\sigma_q=\min\{n\geq0:(J_q)_d=(I^{[q]})_d
                   \text{ for every }d\geq n\},
\]
then $\sigma_q\geq3q/2+c_q+1$.
\end{corollary}
\begin{proof}
Theorem~\ref{thm:LC} gives $H^0_{\mathfrak m}(T_q)\ne0$.
Thus $\mathfrak m\in\Ass T_q$, and localization at $\mathfrak m$ has
depth zero. Since $\sqrt{I^{[q]}}=I\ne\mathfrak m$, this associated
point is embedded. The proof of Proposition~\ref{prop:regularity-reduction}
gives a nonzero component of $J_q/I^{[q]}$ in degree $3q/2+c_q$,
which proves the bound for $\sigma_q$.
\end{proof}

The number $\sigma_q$ bounds degrees in the saturation defect, whereas
$\lambda_q$ measures its annihilation by $\mathfrak m$.
The pencil calculation also bounds $\dim_k L_q$, the generalized
Hilbert--Kunz function of $R/I$ at $q$ \cite{EpsteinYao,DaoSmirnov}.
Under suitable hypotheses, Vraciu~\cite{VraciuGHK} used $(\mathrm{LC})$
to express this function in terms of ordinary Hilbert--Kunz functions.
Brenner and Caminata~\cite{BC} computed its form in graded dimension two.
For the corresponding length problem for ordinary powers, see
Dao and Monta\~no~\cite{DaoMontano}.

\subsection{Associated primes of Frobenius powers}
\label{subsec:ass}
Katzman~\cite{KatzAss} exhibited Frobenius powers whose associated
primes have infinite union. Swanson~\cite{SwaInf} obtained this
phenomenon for a hypersurface domain, and Singh and Swanson~\cite{SinghSwa}
studied it in relation to local cohomology. We determine the set at
each Frobenius power by identifying the parameter values where the fibre
length jumps. Dinh~\cite{DinhAss} studied associated primes in the
point-$S_4$ family of Monsky~\cite{MonskyPoint} used in the example of Brenner and Monsky~\cite{BM}. For the line-$S_4$ family,
the parameters at which the Hilbert--Kunz function jumps give all
height-one associated primes over $k[u,w]$. Together with the vertex,
they account for the full associated-prime set.

\begin{lemma}\label{lem:fibre-torsion}
Let $S$ be a domain, let $N$ be a finite $S$-module of generic rank
$\rho$, and let $\mathfrak p$ be a height-one prime such that
$S_{\mathfrak p}$ is a discrete valuation ring. Then
\[
\mathfrak p\in\Ass_S(N)
\quad\Longleftrightarrow\quad
\dim_{\kappa(\mathfrak p)}(N\otimes_S\kappa(\mathfrak p))>\rho.
\]
\end{lemma}
\begin{proof}
Over the DVR $S_{\mathfrak p}$, write
$N_{\mathfrak p}\simeq S_{\mathfrak p}^{\rho}
\oplus\bigoplus_i S_{\mathfrak p}/(\pi^{a_i})$
with $a_i>0$. The fibre has dimension $\rho$ plus the number of
nonzero torsion summands. These occur precisely when
$\mathfrak pS_{\mathfrak p}$ is an associated prime of
$N_{\mathfrak p}$.
\end{proof}

\begin{theorem}\label{thm:ass-exact}
Let $q=2^e\geq8$ and $r_e=\lfloor(e-1)/2\rfloor$. Let $\mathcal F_e$
be the set of monic irreducible factors in $k[T]$ of
$P_1\cdots P_{r_e}$. Write $g^{\mathrm h}=w^{\deg g}g(u/w)$ and let
$\nu_k(P_j)$ count the distinct monic irreducible factors of $P_j$.
Then
\[
 \Ass_R(R/I^{[q]})=\{I,\mm,(I,w)\}
 \cup\{(I,g^{\mathrm h}):g\in\mathcal F_e\}.
\]
In particular,
\[
\#\Ass_R(R/I^{[q]})
=3+\sum_{j=1}^{r_e}\nu_k(P_j).
\]
If $k$ is algebraically closed, the formula becomes
\[
\#\Ass_R(R/I^{[2^e]})
=3+\frac{8^{\lfloor(e-1)/2\rfloor}-1}{7}.
\]
\end{theorem}
\begin{proof}
Put $S=k[u,w]$ and $T_q=R/I^{[q]}$.
The graded inclusion $S\to T_q$ is finite and $(x,y,z)T_q$ is nilpotent,
with quotient $S$. Thus the unique prime of $T_q$ above a prime
$\mathfrak p\subset S$ is $\mathfrak pT_q+(x,y,z)T_q$.
The contraction formula for associated primes of finite modules under
a finite extension identifies $\Ass_S T_q$ with the contractions of
$\Ass_{T_q}T_q$. Taking inverse images under $R\to T_q$ gives the
corresponding primes of $R$.

The generic $S$-rank of $T_q$ is
\begin{equation}\label{eq:generic-fibre-length}
\rho_q=3q^2-4.
\end{equation}
This follows from \eqref{eq:HK-fibre} at the generic parameter.
At $u=0$, Proposition~\ref{prop:rank} gives surjectivity in target
degree $3q/2+1$ and hence in every higher degree. Socle duality gives
injectivity in the complementary degrees. Thus multiplication by $h$
has maximal rank in every degree, as does generic multiplication by
$h+(u/w)z^4$. The two fibre lengths agree, excluding $(u)$.

Since $P_j$ are squarefree and pairwise coprime
(Proposition~\ref{prop:escape}), Lemma~\ref{lem:fibre-torsion}
shows that the finite nonzero height-one associated primes over $S$
are precisely the homogeneous prime factors of $P_j(u/w)$ for
$1\leq j\leq r_e$. This statement descends from $\bar k$ to $k$:
the fibre dimension over the residue field of an irreducible factor
is unchanged by extending that residue field algebraically.

At the remaining point of the parameter line, $w=0$, the relation
reduces to $z^4$ because $u$ is a unit there. For $q\geq4$ the fibre
has length
\[
\dim_k A_q/(z^4)=4q^2>\rho_q.
\]
Hence $(w)\in\Ass_S T_q$.
The zero prime of $S$ is associated because $T_q$ has positive
rank. By Theorem~\ref{thm:LC},
$H^0_{(u,w)}(T_q)=H^0_{\mathfrak m}(T_q)\ne0$ for $q\geq8$, so $(u,w)\in\Ass_S T_q$.

Every associated prime of a finite graded module over $k[u,w]$
is homogeneous, and every homogeneous prime of $k[u,w]$ is either
$(0)$, a homogeneous height-one prime, or $(u,w)$.
The preceding cases exhaust these possibilities. Lifting them to
$R$ proves the asserted list. Counting the distinct irreducible
factors gives the first cardinality formula. Over an algebraically
closed field, every $P_j$ has $8^{j-1}$ distinct roots. Summing
$\sum_{j=1}^{r_e}8^{j-1}$ gives the second formula.
\end{proof}

\begin{corollary}\label{cor:ass-growth}
For every field $k$ of characteristic two,
\[
\bigcup_{e\geq0}\Ass_R(R/I^{[2^e]})
\quad\text{is infinite},
\]
and these associated-prime sets form an ascending chain.
If $k$ is algebraically closed, then
$\#\Ass_R(R/I^{[2^e]})=\Theta(2^{3e/2})$.
For $q=1,2$ the only associated prime is $I$. For $q=4$ the set is
$\{I,(I,w)\}$.
\end{corollary}
\begin{proof}
The polynomials $P_j$ are pairwise coprime of positive degree,
so they contribute a new irreducible factor at every new escape
level, over any base field. The theorem and its formula for $r_e$
give the inclusions and, when $k$ is algebraically closed, the
growth assertion. For $q=1,2$ the relation $uz^4+wh$ vanishes in
$A_q[u,w]$, so $T_q$ is free over $S$. For $q=4$ we have
$T_4=A_4[u,w]/(wh)$. The variable $u$ is a nonzerodivisor,
and the only height-one torsion prime over $S$ is $(w)$.
\end{proof}

\begin{example}
At $q=32$ there are twelve associated primes over $\overline{\mathbb F}_2$. Over $\mathbb F_2$,
$P_1(T)=T+1$ and
\[
P_2(T)=(T^2+T+1)(T^6+T^5+T^3+T^2+1),
\]
so $\#\Ass_R(R/I^{[32]})=6$.
\end{example}

\subsection{Embedded primary components}
\begin{theorem}\label{thm:primary}
For every dyadic $q\geq8$, each reduced primary decomposition of
$I^{[q]}$ contains an $\mm$-primary component $Q$ satisfying
\[
 w^{c_q-1}\notin Q.
\]
Consequently $\mm^{N_q}\subseteq Q$ requires $N_q\geq c_q$, and
no fixed $\ell$ satisfies~\eqref{eq:SS-question} for all $q$.
\end{theorem}
\begin{proof}
By Theorem~\ref{thm:LC}, choose $a\in R$ whose class in $L_q$
satisfies $w^{c_q-1}a\notin I^{[q]}$. Since $L_q\ne0$, the maximal
ideal $\mm$ is associated, so its component $Q$ occurs.
For any other component $Q'$ choose $s\in\mm\setminus\sqrt{Q'}$.
Some power of $s$ kills the class of $a$, whence $s^ta\in Q'$
for some $t$. Since $Q'$ is primary, we have $a\in Q'$. If $w^{c_q-1}\in Q$,
then $w^{c_q-1}a$ belongs to every component, a contradiction.

If $\mm^{N_q}\subseteq Q$ then $w^{N_q}\in Q$, hence $N_q\geq c_q$.
Likewise $(\mm^{[q]})^\ell\subseteq Q$ implies $q\ell\geq c_q$.
Proposition~\ref{prop:index-lower} gives $c_q/q\to\infty$.
Removing redundant components and intersecting those with the same
radical preserves the containments in~\eqref{eq:SS-question}.
\end{proof}

\begin{corollary}\label{cor:ordinary-primary}
There is no constant $\gamma$ such that every $I^{[q]}$ has a primary
decomposition with
$P_{i,q}^{\lceil\gamma q\rceil}\subseteq Q_{i,q}$ for every component.
For $q=4^r$, $r\geq2$, every $\mm$-primary component in a reduced
primary decomposition requires an ordinary radical-power containment
exponent at least $q^{3/2}/8-1$.
The same conclusions hold for $J+K^{[q]}\subseteq P$ from
Corollary~\ref{cor:polynomial-regularity}, with maximal ideal $P_+$.
\end{corollary}
\begin{proof}
Theorem~\ref{thm:primary} bounds every such exponent below by $c_q$.
Use Proposition~\ref{prop:index-lower} and the equality
$P/(J+K^{[q]})=R/I^{[q]}$.
\end{proof}

\subsection{Gr\"obner bases}
The same local-cohomology class also bounds the degree in the least variable.

\begin{lemma}\label{lem:grevlex-saturation}
Let $F\subseteq k[X_1,\ldots,X_n]$ be a homogeneous ideal. Fix graded
reverse lexicographic order with $X_n$ the least variable. If $D$ is the maximum degree in $X_n$ among its reduced
Gr\"obner-basis elements, then
\[
X_n^D(F:X_n^\infty)\subseteq F.
\]
\end{lemma}
\begin{proof}
Write each basis element as $g=X_n^{a_g}\widetilde g$, where
$X_n\nmid\ini(\widetilde g)$. This is possible because,
for a homogeneous polynomial in this order, the exponent of $X_n$ in
its leading monomial is the smallest exponent occurring in any term.
The polynomials $\widetilde g$ form a Gr\"obner basis of
$F:X_n^\infty$: if $X_n^bf\in F$, some
$\ini(g)$ divides $X_n^b\ini(f)$, so
$\ini(\widetilde g)$ divides $\ini(f)$.
Repeated division proves the assertion. Since $a_g\leq D$, each
$X_n^D\widetilde g$ belongs to $F$, giving the required containment.
\end{proof}

\begin{theorem}\label{thm:groebner}
Let $\mathfrak L_q=J+K^{[q]}\subseteq P=k[x,y,z,u,w]$ as in
Corollary~\ref{cor:polynomial-regularity}. Fix graded reverse
lexicographic order with $w$ the least variable. Let $\mathcal G_q$ be the reduced Gr\"obner basis of $\mathfrak L_q$
and put $D_q=\max_{g\in\mathcal G_q}\deg_w g$. Then
\[
D_q\geq c_q.
\]
In particular, $D_q/q\to\infty$ along powers of two.
\end{theorem}
\begin{proof}
Lemma~\ref{lem:grevlex-saturation} shows that $w^{D_q}$ annihilates
$H^0_{P_+}(P/\mathfrak L_q)$, since every class in this module is killed
by a power of $w$. This module equals $L_q$. Theorem~\ref{thm:LC}
gives $w^{c_q-1}L_q\ne0$, so $D_q\geq c_q$.
\end{proof}

\begin{thebibliography}{99}
\bibitem{Asgharzadeh}
M. Asgharzadeh, \emph{On the (LC) conjecture}, preprint.
\href{https://arxiv.org/abs/1512.02518}{arXiv:1512.02518}.
\bibitem{Borevitz}
L. Borevitz, N. Nader, T. J. Sandstrom, A. Shapiro, A. Simpson, and
J. Zomback, \emph{On localization of tight closure in line-$S_4$
quartics}, J. Pure Appl. Algebra \textbf{228} (2024), no.~9, 107682.
\bibitem{Brenner}
H. Brenner, \emph{A linear bound for Frobenius powers and an inclusion
bound for tight closure}, Michigan Math. J. \textbf{53} (2005), no.~3,
585--596.
\bibitem{BrennerStable}
H. Brenner, \emph{Looking out for stable syzygy bundles}, Adv. Math.
\textbf{219} (2008), no.~2, 401--427.
\bibitem{BC}
H. Brenner and A. Caminata, \emph{Generalized Hilbert--Kunz function
in graded dimension two}, Nagoya Math. J. \textbf{230} (2018), 1--17.
\bibitem{BM}
H. Brenner and P. Monsky, \emph{Tight closure does not commute with
localization}, Ann. of Math. (2) \textbf{171} (2010), no.~1, 571--588.
\bibitem{Chardin}
M. Chardin, \emph{On the behavior of Castelnuovo--Mumford regularity
with respect to some functors}, preprint.
\href{https://arxiv.org/abs/0706.2731}{arXiv:0706.2731}.
\bibitem{ChardinSurvey}
M. Chardin, \emph{Some results and questions on Castelnuovo--Mumford
regularity}, in: Syzygies and Hilbert functions, Lect. Notes Pure
Appl. Math. \textbf{254}, Chapman \& Hall/CRC, Boca Raton, FL, 2007,
1--40.
\bibitem{DaoMontano}
H. Dao and J. Monta\~no, \emph{Length of local cohomology of powers
of ideals}, Trans. Amer. Math. Soc. \textbf{371} (2019), no.~5,
3483--3503.
\bibitem{DaoSmirnov}
H. Dao and I. Smirnov, \emph{On generalized Hilbert--Kunz function
and multiplicity}, Israel J. Math. \textbf{237} (2020), no.~1,
155--184. \href{https://arxiv.org/abs/1305.1833}{arXiv:1305.1833}.
\bibitem{DSNP}
A. De Stefani, L. N\'u\~nez-Betancourt, and F. P\'erez,
\emph{On the existence of $F$-thresholds and related limits},
Trans. Amer. Math. Soc. \textbf{370} (2018), no.~9, 6629--6650.
\bibitem{Dinh}
T. T. Dinh, \emph{Growth of primary decompositions of Frobenius powers
of ideals}, J. Algebra \textbf{321} (2009), no.~3, 829--846.
\bibitem{DinhAss}
T. T. Dinh, \emph{Associated primes of the example of Brenner and
Monsky}, Comm. Algebra \textbf{41} (2013), no.~1, 109--123.
\bibitem{EpsteinYao}
N. Epstein and Y. Yao, \emph{Some extensions of Hilbert--Kunz
multiplicity}, Collect. Math. \textbf{68} (2017), 69--85.
\href{https://arxiv.org/abs/1103.4730}{arXiv:1103.4730}.
\bibitem{HS}
S. Hermiller and I. Swanson, \emph{Computations with Frobenius powers},
Experiment. Math. \textbf{14} (2005), no.~2, 161--173.
\bibitem{HH}
M. Hochster and C. Huneke, \emph{Tight closure, invariant theory, and
the Brian\c{c}on--Skoda theorem}, J. Amer. Math. Soc. \textbf{3}
(1990), no.~1, 31--116.
\bibitem{HunekeLC}
C. Huneke, \emph{The saturation of Frobenius powers of ideals},
Comm. Algebra \textbf{28} (2000), no.~12, 5563--5572.
\bibitem{Katzman}
M. Katzman, \emph{The complexity of Frobenius powers of ideals},
J. Algebra \textbf{203} (1998), no.~1, 211--225.
\bibitem{KatzAss}
M. Katzman, \emph{Finiteness of $\bigcup_e\Ass F^e(M)$ and its
connections to tight closure}, preprint.
\href{https://arxiv.org/abs/math/0209355}{arXiv:math/0209355}.
\bibitem{KSSZ}
M. Katzman, K. Schwede, A. K. Singh, and W. Zhang, \emph{Rings of
Frobenius operators}, Math. Proc. Cambridge Philos. Soc. \textbf{157}
(2014), no.~1, 151--167.
\bibitem{KZ}
M. Katzman and W. Zhang, \emph{Castelnuovo--Mumford regularity and the
discreteness of $F$-jumping coefficients in graded rings},
Trans. Amer. Math. Soc. \textbf{366} (2014), no.~7, 3519--3533.
\bibitem{MonskyPoint}
P. Monsky, \emph{Hilbert--Kunz functions in a family: point-$S_4$
quartics}, J. Algebra \textbf{208} (1998), no.~1, 343--358.
\bibitem{Monsky}
P. Monsky, \emph{Hilbert--Kunz functions in a family: line-$S_4$
quartics}, J. Algebra \textbf{208} (1998), no.~1, 359--371.
\bibitem{PQ}
T. Polstra and P. H. Quy, \emph{Nilpotence of Frobenius actions on
local cohomology and Frobenius closure of ideals},
J. Algebra \textbf{529} (2019), 196--225.
\bibitem{SN}
R. Y. Sharp and N. Nossem, \emph{Ideals in a perfect closure, linear
growth of primary decompositions, and tight closure},
Trans. Amer. Math. Soc. \textbf{356} (2004), no.~9, 3687--3720.
\bibitem{SinghSwa}
A. K. Singh and I. Swanson, \emph{Associated primes of local
cohomology modules and of Frobenius powers}, Int. Math. Res. Not.
(2004), no.~33, 1703--1733.
\bibitem{SS}
K. E. Smith and I. Swanson, \emph{Linear bounds on growth of associated
primes for monomial ideals}, Comm. Algebra \textbf{25} (1997), no.~10,
3071--3079.
\bibitem{Swanson}
I. Swanson, \emph{Powers of ideals: primary decompositions, Artin--Rees
lemma and regularity}, Math. Ann. \textbf{307} (1997), no.~2, 299--313.
\bibitem{SwaInf}
I. Swanson, \emph{Infinitely many associated primes of Frobenius
powers and local cohomology}, manuscript.
\url{https://www.math.purdue.edu/~iswanso/frobprim.pdf}.
\bibitem{Trivedi}
V. Trivedi, \emph{Semistability and Hilbert--Kunz multiplicities for
curves}, J. Algebra \textbf{284} (2005), no.~2, 627--644.
\bibitem{Vraciu}
A. Vraciu, \emph{Local cohomology of Frobenius images over graded
affine algebras}, J. Algebra \textbf{228} (2000), no.~1, 347--356.
\bibitem{VraciuGHK}
A. Vraciu, \emph{An observation on generalized Hilbert--Kunz
functions}, Proc. Amer. Math. Soc. \textbf{144} (2016), no.~8,
3221--3229. \href{https://arxiv.org/abs/1510.00668}{arXiv:1510.00668}.
\bibitem{Yao}
Y. Yao, \emph{Primary decomposition: compatibility, independence and
linear growth}, Proc. Amer. Math. Soc. \textbf{130} (2002), no.~6,
1629--1637.
\end{thebibliography}
\end{document}
